%% ============================================================================
%%  CAPITOLO 5 — Derivate
%%  Fonte principale per argomenti e dimostrazioni: Marcellini–Sbordone, cap. 5
%%  (§39 tasso di accrescimento, §40 definizione di derivata, §41 operazioni,
%%  §42 funzioni composte e inverse, §43 derivate delle funzioni elementari,
%%  §44 significato geometrico, retta tangente; appendice §45 funzioni
%%  trigonometriche inverse).
%% ============================================================================
\aperturacapitolo{Derivate}{cap:derivate}

\epigraph{Nova methodus pro maximis et minimis, itemque tangentibus, quae nec fractas
nec irrationales quantitates moratur, et singulare pro illis calculi genus.\\[2pt]
\upshape(Un nuovo metodo per i massimi e i minimi, e anche per le tangenti, che non si
arresta davanti alle quantità fratte o irrazionali, e un singolare genere di calcolo per
esse.)}{Gottfried Wilhelm Leibniz, titolo della memoria negli \emph{Acta Eruditorum}
(1684)}

\lettrine[lines=3, findent=2pt, nindent=0pt]{C}{on la derivata} l'Analisi diventa uno
strumento di calcolo. I capitoli precedenti ci hanno dato il linguaggio dei limiti e le
proprietà delle funzioni continue; ora usiamo quel linguaggio per rispondere a una
domanda che attraversa tutta la fisica e l'ingegneria: \emph{con che rapidità varia una
grandezza?} La velocità di un velivolo è la rapidità con cui varia la sua posizione,
il rateo di salita è la rapidità con cui varia la quota, il consumo è la rapidità con
cui diminuisce la massa di combustibile, il gradiente termico è la rapidità con cui la
temperatura cambia con la quota. In tutti questi casi si confronta la variazione della
grandezza con la variazione della variabile da cui dipende, in un intervallo sempre più
piccolo: si calcola cioè il limite di un \emph{rapporto incrementale}.

La stessa operazione ha un significato geometrico altrettanto importante: la derivata è
la pendenza della retta tangente al grafico. Vedremo che una funzione derivabile, vicino
a un punto, «assomiglia» alla sua retta tangente: è l'idea di \emph{linearizzazione},
che sta alla base di gran parte dei modelli ingegneristici, dalle piccole perturbazioni
in meccanica del volo all'analisi della propagazione degli errori di misura. Seguiremo
il testo di riferimento \cite{marcellini-sbordone}: definizione, regole di calcolo,
derivate delle funzioni elementari, interpretazione geometrica; in appendice
introdurremo le funzioni trigonometriche inverse, indispensabili per esprimere gli
angoli a partire dai loro seni, coseni e tangenti.

\medskip
\noindent\textbf{In questo capitolo.}
\begin{itemize}
  \item Il tasso di accrescimento e la velocità istantanea (\S\ref{sec:tasso}).
  \item La definizione di derivata; derivabilità e continuità; derivate successive
  (\S\ref{sec:def-derivata}).
  \item Le regole di derivazione: somma, prodotto, quoziente (\S\ref{sec:operazioni-derivate}),
  funzioni composte e funzioni inverse (\S\ref{sec:composte-inverse}).
  \item Le derivate delle funzioni elementari (\S\ref{sec:derivate-elementari}).
  \item La retta tangente e l'approssimazione lineare; le funzioni $x^{n}\sen\frac1x$
  (\S\ref{sec:retta-tangente}).
  \item Appendice: le funzioni trigonometriche inverse (\S\ref{sec:trig-inverse}).
  \item Esercizi di riepilogo (\S\ref{sec:esercizi-cap5}).
  \item Punti fondamentali del capitolo (\S\ref{sec:sintesi-cap05}).
\end{itemize}

%% ----------------------------------------------------------------------------
%%  §5.1  Tasso di accrescimento. Significato meccanico della derivata
%% ----------------------------------------------------------------------------
\section{Tasso di accrescimento. Significato meccanico della derivata}
\label{sec:tasso}

Consideriamo un semplice processo di crescita di un corpo, supponendo che il peso $p$
del corpo vari al crescere del tempo; cioè supponiamo che il peso sia una funzione del
tempo, $p=p(t)$. Prendiamo in considerazione le variazioni di peso a partire da un certo
istante $t$. Al tempo $t$ il peso è $p(t)$, mentre all'istante $t+h$, dopo che è
trascorso un tempo uguale ad $h$, il peso è $p(t+h)$. Quindi, nell'intervallo di tempo
$h$, il cambio di peso è $p(t+h)-p(t)$. Il rapporto
\begin{equation}\label{eq:tasso-medio}
  \frac{p(t+h)-p(t)}{h}
\end{equation}
dà un'indicazione di quanto sia cambiato il peso per unità di tempo. Più precisamente
il rapporto \eqref{eq:tasso-medio} esprime la variazione \emph{media} per unità di tempo
del peso nell'intervallo $[t,t+h]$ (o $[t+h,t]$ se $h<0$). Si chiama anche \emph{tasso
medio di accrescimento}, o tasso medio di variazione, o \emph{velocità media di
accrescimento}.

Invece della variazione media nell'intervallo $[t,t+h]$, spesso è più utile la
variazione \emph{istantanea} al tempo $t$. Intuitivamente si considera l'espressione
\eqref{eq:tasso-medio} per alcuni valori di $h$ sempre più vicini a zero. Con il
linguaggio più preciso introdotto nei capitoli precedenti, si calcola il limite, per
$h\to0$, del rapporto \eqref{eq:tasso-medio}:
\begin{equation}\label{eq:tasso}
  \text{tasso di accrescimento}=\lim_{h\to0}\frac{p(t+h)-p(t)}{h}.
\end{equation}
Si noti come sia indispensabile usare il limite per $h\to0$, dato che, se nella
\eqref{eq:tasso-medio} si pone direttamente $h=0$, si ottiene un'espressione senza
significato. Invece, con il linguaggio dei limiti, il limite in \eqref{eq:tasso} è una
forma indeterminata $\frac00$ (se $p$ è una funzione continua), del tipo già incontrato
nell'esempio della velocità istantanea del \cref{sec:premessa-funzioni}.

Per mostrare come il tasso di accrescimento possa essere effettivamente calcolato,
consideriamo a titolo di esercizio il caso in cui il peso dipenda dal tempo in modo
quadratico,
\begin{equation}\label{eq:peso-quadratico}
  p(t)=t^{2};
\end{equation}
per ogni $h\neq0$ risulta
\begin{equation}\label{eq:tasso-t2}
  \frac{p(t+h)-p(t)}{h}=\frac{(t+h)^{2}-t^{2}}{h}=\frac{t^{2}+2ht+h^{2}-t^{2}}{h}=2t+h .
\end{equation}
L'ultimo membro della \eqref{eq:tasso-t2} tende a $2t$ per $h\to0$. Quindi il tasso di
accrescimento vale $2t$. Ciò significa che, al crescere del tempo $t>0$, non soltanto il
peso cresce come $t^{2}$, ma anche il \emph{cambiamento} di peso per unità di tempo
aumenta (nel caso in considerazione, in modo proporzionale al tempo).

Proponiamo un esempio numerico: secondo la legge $p(t)=t^{2}$, al tempo $t=10$ il peso
risulta essere uguale a $p(10)=100$. Il tasso di accrescimento, uguale a $2t$, al tempo
$t=10$ vale $20$. Ciò significa che, dopo una unità di tempo, il peso del corpo aumenta
di circa $20$ unità; quindi $p(11)$ vale all'incirca $120$. Si noti che effettivamente il
valore trovato, $120$, non differisce di molto da $p(11)=11^{2}=121$; approfondiremo
questo aspetto nel \cref{sec:retta-tangente} e, in modo più completo, con la formula di
Taylor. Il programma seguente calcola i tassi medi per $h$ sempre più piccoli, in
aritmetica razionale esatta.

\begin{python}[Tassi medi di accrescimento per $p(t)=t^{2}$ e $t=10$]
from fractions import Fraction as F

p = lambda t: t**2                    # peso in funzione del tempo
t = F(10)
for k in range(0, 7):
    h = F(1, 10**k)                   # aritmetica razionale esatta
    tasso_medio = (p(t + h) - p(t)) / h
    print(f"h = {str(h):>9}   tasso medio = {float(tasso_medio)}")
print("previsione lineare p(11) ~", p(t) + 2 * t * 1,
      "  valore vero:", p(t + 1))
\end{python}
\begin{risultato}
h =         1   tasso medio = 21.0
h =      1/10   tasso medio = 20.1
h =     1/100   tasso medio = 20.01
h =    1/1000   tasso medio = 20.001
h =   1/10000   tasso medio = 20.0001
h =  1/100000   tasso medio = 20.00001
h = 1/1000000   tasso medio = 20.000001
previsione lineare p(11) ~ 120   valore vero: 121
\end{risultato}

I tassi medi valgono esattamente $20+h$, come previsto dalla \eqref{eq:tasso-t2}, e si
avvicinano a $20$; ma, ancora una volta, il calcolo con alcuni valori di $h$ suggerisce
il limite, mentre è la \eqref{eq:tasso-t2} a dimostrarlo.

\subsection{La velocità istantanea}

Abbiamo già detto che «velocità di accrescimento» è sinonimo di «tasso di
accrescimento»; ciò deriva dal fatto che una velocità si definisce in modo analogo a
quanto fatto sopra. Consideriamo ad esempio un velivolo che percorre una rotta
rettilinea, e indichiamo con $s(t)$ lo spazio percorso in funzione del tempo. La
\emph{velocità media} nell'intervallo di tempo $[t,t+h]$ è uguale al rapporto tra lo
spazio percorso $s(t+h)-s(t)$ e il tempo $h$ impiegato a percorrerlo. La \emph{velocità
istantanea} (quella che, a meno del vento, si ricava dagli strumenti di bordo) è il
limite, per $h\to0$, della velocità media; quindi
\begin{equation}\label{eq:velocita-istantanea}
  \text{velocità istantanea}=\lim_{h\to0}\frac{s(t+h)-s(t)}{h}.
\end{equation}
Il significato geometrico è illustrato nella \cref{fig:velocita-media}: la velocità media
è la pendenza della retta che congiunge due punti del grafico della legge oraria; la
velocità istantanea è la pendenza limite di tali rette quando il secondo punto si
avvicina al primo.

\begin{figure}[htb]
\centering
\begin{tikzpicture}
\begin{axis}[width=0.9\textwidth, height=5.6cm, xmin=0, xmax=6, ymin=0, ymax=5.6,
  axis lines=left, xtick={2,4.6}, xticklabels={$t$,$t+h$}, ytick={1.404,4.6221},
  yticklabels={$s(t)$,$s(t+h)$}, xlabel={tempo}, ylabel={spazio}, clip=false]
  \addplot[curva, domain=0.2:5.2, samples=100] {0.5+0.25*x^2-0.012*x^3};
  \addplot[tratteggio] coordinates {(0,1.404) (2,1.404) (2,0)};
  \addplot[tratteggio] coordinates {(0,4.6221) (4.6,4.6221) (4.6,0)};
  \addplot[curva2, domain=1.2:5.1] {1.404+(4.6221-1.404)/2.6*(x-2)};
  \addplot[curva3, domain=0.4:4.2] {1.404+0.856*(x-2)};
  \addplot[only marks, mark=*, mark size=1.6pt, blu] coordinates {(2,1.404) (4.6,4.6221)};
  \node[font=\scriptsize, rame, anchor=west] at (axis cs:5.1,5.0) {secante};
  \node[font=\scriptsize, verde!50!black, anchor=north west] at (axis cs:3.6,2.35) {tangente};
\end{axis}
\end{tikzpicture}
\caption{La velocità media nell'intervallo $[t,t+h]$ è la pendenza della retta secante
(rame) per i punti $(t,s(t))$ e $(t+h,s(t+h))$ del grafico della legge oraria; quando
$h\to0$ la secante tende alla retta tangente (verde), la cui pendenza è la velocità
istantanea.}
\label{fig:velocita-media}
\end{figure}

È chiaro che nei due esempi precedenti lo schema matematico è identico. In entrambi gli
esempi occorre calcolare \emph{il limite di un rapporto incrementale}, così chiamato
perché a denominatore c'è l'incremento $h$ della variabile indipendente, mentre a
numeratore c'è il corrispondente incremento della variabile dipendente. Occorre
calcolare il limite del rapporto incrementale anche in molte altre situazioni, analoghe
a quelle dei due esempi esposti:
\begin{itemize}
  \item l'\emph{accelerazione} di un corpo in moto rettilineo è il limite del rapporto
  incrementale della velocità rispetto al tempo;
  \item la \emph{portata in massa} di combustibile consumato da un motore è il limite del
  rapporto incrementale della massa di combustibile bruciata rispetto al tempo;
  \item la \emph{densità lineare} di una trave non omogenea è il limite del rapporto tra
  la massa di un tratto di trave e la sua lunghezza, quando la lunghezza tende a zero;
  \item il \emph{gradiente termico} verticale dell'atmosfera è il limite del rapporto
  incrementale della temperatura rispetto alla quota.
\end{itemize}
Un altro esempio, di tipo geometrico, è studiato nel \cref{sec:retta-tangente}.
Introdurremo nel prossimo paragrafo la \emph{derivata} come limite del rapporto
incrementale, quando l'incremento tende a zero.

\begin{idea}[Il variometro: uno strumento che misura una derivata]
Il variometro, presente sul cruscotto di ogni aliante e di quasi ogni velivolo, indica
la \emph{velocità verticale}, cioè il tasso di variazione della quota rispetto al tempo,
$\lim_{h\to0}\frac{z(t+h)-z(t)}{h}$. Il variometro meccanico realizza fisicamente
un rapporto incrementale: una capsula, collegata direttamente alla presa statica, è
racchiusa in un serbatoio che comunica con la stessa presa soltanto attraverso un
sottile tubo capillare, che rallenta l'equalizzazione delle pressioni. Quando la quota
cambia, la pressione nella capsula segue subito quella esterna, mentre quella nel
serbatoio resta indietro, e la differenza di pressione che si stabilisce è, con buona
approssimazione, proporzionale alla rapidità di variazione della pressione statica, e
quindi della quota. Lo strumento ha un ritardo, perché in effetti confronta valori in
istanti diversi: misura cioè un rapporto incrementale con un incremento di tempo
piccolo ma non nullo. I variometri elettronici calcolano il rapporto incrementale per
via numerica, con gli accorgimenti che vedremo nel \cref{sec:def-derivata}.
\end{idea}

\begin{notastorica}[Newton, Leibniz e le notazioni]
Il calcolo differenziale nacque due volte, in modo indipendente. Isaac Newton lo
sviluppò negli anni 1665--1666 come «metodo delle flussioni». Per lui le grandezze sono
quantità che «fluiscono» nel tempo, le \emph{fluenti}; la loro velocità di variazione è
la \emph{flussione}, e il prodotto della flussione per un intervallo di tempo
infinitamente piccolo $o$ è il \emph{momento}, cioè l'incremento infinitesimo della
fluente. Il problema centrale, nel trattato \emph{De methodis serierum et fluxionum}
scritto nel 1671, è proprio il nostro: data una relazione tra quantità fluenti, trovare
la relazione tra le loro flussioni. Poiché la variabile indipendente è sempre il tempo,
la derivata di $y$ rispetto a $x$ compare come rapporto tra le flussioni di $y$ e di $x$.
Il trattato fu pubblicato soltanto nel 1736, in traduzione inglese (\emph{The Method of
Fluxions and Infinite Series}); la notazione col punto, $\dot x$, ancora oggi d'uso in
meccanica per le derivate rispetto al tempo, Newton la adottò soltanto alla fine del 1691.
Gottfried Wilhelm Leibniz pubblicò per primo, nell'ottobre 1684, la memoria citata
nell'epigrafe, uscita a Lipsia negli \emph{Acta Eruditorum}: poche pagine che rendevano
pubblici i differenziali $dx$, $dy$, ideati nei manoscritti parigini del 1675, e le regole
per calcolarli. Leibniz pensava le grandezze come formate da parti infinitesime, e la
sua notazione, che conserva nel simbolo $\frac{dy}{dx}$ l'origine della derivata come
rapporto di incrementi, si è dimostrata straordinariamente efficace, come vedremo con la
regola della catena. Sulla disputa per la priorità che seguì si veda la nota storica
del \S\ref{sec:teorema-fondamentale}; il nome «derivata» e il simbolo $f'(x)$ sono più
tardi (si veda la nota storica alla fine del \S\ref{sec:def-derivata}).
\end{notastorica}

%% ----------------------------------------------------------------------------
%%  §5.2  Definizione di derivata
%% ----------------------------------------------------------------------------
\section{Definizione di derivata}
\label{sec:def-derivata}

\begin{definizione}[Derivata]\label{def:derivata}\index{derivata}\index{funzione!derivabile}
Sia $f$ definita nell'intervallo aperto $(a,b)$ e sia $x$ un punto di $(a,b)$. Si dice
che la funzione $f$ è \emph{derivabile} nel punto $x$ se \emph{esiste finito il limite
del rapporto incrementale}
\begin{equation}\label{eq:def-derivata}
  \lim_{h\to0}\frac{f(x+h)-f(x)}{h}.
\end{equation}
Tale limite è la \emph{derivata} di $f$ nel punto $x$, e si indica con una delle
seguenti notazioni, fra loro equivalenti:
\begin{equation}\label{eq:notazioni-derivata}
  f'(x), \qquad \frac{df}{dx}(x), \qquad Df(x), \qquad y', \qquad \frac{dy}{dx}, \qquad Dy .
\end{equation}
Si dice che $f$ è \emph{derivabile nell'intervallo aperto} $(a,b)$ se è derivabile in ogni
punto $x\in(a,b)$.
\end{definizione}

Qualche precisazione sulla definizione. Il rapporto incrementale è una funzione della
sola variabile $h$, definita per i valori $h\neq0$ tali che $x+h\in(a,b)$; poiché $x$ è
interno ad $(a,b)$, questo insieme contiene un intorno di $0$ privato di $0$, e il limite
per $h\to0$ ha senso (\cref{def:limite-funzione}). Le notazioni di Leibniz
$\frac{df}{dx}$, $\frac{dy}{dx}$ ricordano che la derivata è il limite di un rapporto tra
l'incremento $\Delta y=f(x+h)-f(x)$ della variabile dipendente e l'incremento
$\Delta x=h$ della variabile indipendente; non sono, per ora, quozienti di numeri. Con il
cambiamento di variabile $\xi=x+h$ si ottiene la forma equivalente, spesso comoda,
\begin{equation}\label{eq:derivata-xi}
  f'(x)=\lim_{\xi\to x}\frac{f(\xi)-f(x)}{\xi-x}
\end{equation}
(le due forme sono equivalenti per il \cref{teo:limite-composta}, applicato con la
funzione interna $g(\xi)=\xi-x$, diversa da $0$ per $\xi\neq x$, e viceversa con
$g(h)=x+h$, diversa da $x$ per $h\neq0$).
Infine, la derivata è una nozione \emph{locale}: se due funzioni coincidono in un intorno
di $x$, i loro rapporti incrementali coincidono per $h$ piccolo, e per la località del
limite (\cref{lem:localita}) esse hanno la stessa derivata in $x$, oppure nessuna delle
due è derivabile.

In alcuni casi è utile considerare, al posto del limite completo per $h\to0$ nella
\eqref{eq:def-derivata}, soltanto il limite destro per $h\to0^{+}$, oppure il limite per
$h\to0^{-}$. Nel primo caso si parla di \emph{derivata destra}, indicata con
$f'_{+}(x)$, nel secondo caso di \emph{derivata sinistra}, indicata con $f'_{-}(x)$. Per
la \cref{prop:dx-sx}, \emph{$f$ è derivabile in un punto interno $x$ se e solo se
esistono finite e uguali la derivata destra e la derivata sinistra}, e in tal caso
$f'(x)=f'_{+}(x)=f'_{-}(x)$. Se $f$ è definita in $[a,b]$, si dice che $f$ è
\emph{derivabile nell'intervallo chiuso} $[a,b]$ se è derivabile in ogni punto $x\in(a,b)$
e inoltre se $f$ ammette derivata destra (finita) nel punto $x=a$ e derivata sinistra (finita) nel punto
$x=b$.

\subsection{I primi esempi}

\begin{esempio}[Funzioni costanti e lineari]\label{es:derivata-lineare}
Iniziamo dalla funzione costante $f(x)=q$, per ogni $x\in\R$, e proviamo che tale
funzione è derivabile su tutto $\R$ e che la derivata è identicamente nulla; infatti il
rapporto incrementale vale costantemente zero, qualunque sia l'incremento $h\neq0$:
\begin{equation}\label{eq:derivata-costante}
  \frac{f(x+h)-f(x)}{h}=\frac{q-q}{h}=\frac0h=0,
\end{equation}
e quindi anche il limite del rapporto incrementale, per $h\to0$, vale zero (il lettore non
cada nell'errore di considerare il limite per $h\to0$ della \eqref{eq:derivata-costante}
una forma indeterminata $\frac00$: il rapporto è la funzione costante $0$, e il suo limite
è $0$). Più in generale verifichiamo che la derivata della funzione $f(x)=mx+q$, con $m$
e $q$ costanti (il cui grafico è una retta), è identicamente uguale a $m$; infatti il
rapporto incrementale vale costantemente $m$, qualunque sia $h\neq0$:
\begin{equation}\label{eq:derivata-retta}
  \frac{f(x+h)-f(x)}{h}=\frac{[m(x+h)+q]-[mx+q]}{h}=\frac{mh}{h}=m .
\end{equation}
Abbiamo già calcolato nel paragrafo precedente la derivata della funzione $f(x)=x^{2}$,
trovando $f'(x)=2x$ \eqref{eq:tasso-t2}.
\end{esempio}

\begin{esempio}[Il valore assoluto non è derivabile in $0$]\label{es:modulo-non-derivabile}
Verifichiamo invece che la funzione $f(x)=\abs x$ non è derivabile per $x=0$. Infatti, se
$h\neq0$, si ha
\begin{equation}\label{eq:rapporto-modulo}
  \frac{f(0+h)-f(0)}{h}=\frac{\abs{0+h}-\abs0}{h}=\frac{\abs h}{h}.
\end{equation}
Abbiamo già incontrato questa funzione nella \eqref{eq:segno}. Il limite per $h\to0$ del
rapporto incrementale non esiste, perché risulta \eqref{eq:modulo-su-x}
\begin{equation}\label{eq:derivate-modulo}
  \lim_{h\to0^{+}}\frac{\abs h}{h}=1, \qquad \lim_{h\to0^{-}}\frac{\abs h}{h}=-1 .
\end{equation}
Quindi $f(x)=\abs x$ non è derivabile per $x=0$, mentre esistono le derivate destra e
sinistra, uguali rispettivamente a $+1$ e $-1$. Nei punti $x\neq0$ la funzione coincide
localmente con $x$ o con $-x$, ed è quindi derivabile con derivata $1$ per $x>0$ e $-1$
per $x<0$. Geometricamente, il grafico ha in $0$ uno \emph{spigolo}: le due semirette
hanno pendenze diverse (\cref{fig:non-derivabili}, a sinistra).
\end{esempio}

Un altro modo in cui la derivabilità può venire meno è che il rapporto incrementale
tenda all'infinito. Se il limite \eqref{eq:def-derivata} vale $+\infty$ o $-\infty$, si
dice talvolta che $f$ ha in $x$ \emph{derivata infinita}; ma la funzione \emph{non} è
derivabile, perché la definizione richiede un limite finito.

\begin{esempio}[Derivate infinite]\label{es:derivate-infinite}
Per $f(x)=\sqrt[3]x$ (la radice cubica, definita su tutto $\R$ come inversa della
funzione $x^{3}$) si ha, per $h\neq0$,
\[
  \frac{\sqrt[3]h-0}{h}=\frac1{\sqrt[3]{h^{2}}}\to+\infty \qquad (h\to0),
\]
e il grafico ha nell'origine una tangente verticale. Per $f(x)=\sqrt{\abs x}$ si ha invece
$\frac{\sqrt{\abs h}}{h}\to+\infty$ per $h\to0^{+}$ e $\to-\infty$ per $h\to0^{-}$: il
grafico presenta una \emph{cuspide} (\cref{fig:non-derivabili}). In entrambi i casi la
funzione è continua in $0$ ma non derivabile.
\end{esempio}

\begin{figure}[htb]
\centering
\begin{tikzpicture}
\begin{axis}[name=a1, width=0.36\textwidth, height=0.32\textwidth, axis lines=middle,
  xmin=-1.3, xmax=1.3, ymin=-1.2, ymax=1.3, ticks=none, title={\scriptsize $\abs x$: spigolo}]
  \addplot[curva, domain=-1.2:1.2, samples=3] {abs(x)};
\end{axis}
\begin{axis}[at={(a1.east)}, anchor=west, xshift=0.6cm, name=a2, width=0.36\textwidth,
  height=0.32\textwidth, axis lines=middle, xmin=-1.3, xmax=1.3, ymin=-1.2, ymax=1.3,
  ticks=none, title={\scriptsize $\sqrt[3]x$: tangente verticale}]
  \addplot[curva, domain=0:1.2, samples=80] {x^(1/3)};
  \addplot[curva, domain=-1.2:0, samples=80] {-(-x)^(1/3)};
\end{axis}
\begin{axis}[at={(a2.east)}, anchor=west, xshift=0.6cm, width=0.36\textwidth,
  height=0.32\textwidth, axis lines=middle, xmin=-1.3, xmax=1.3, ymin=-1.2, ymax=1.3,
  ticks=none, title={\scriptsize $\sqrt{\abs x}$: cuspide}]
  \addplot[curva, domain=-1.2:1.2, samples=161] {sqrt(abs(x))};
\end{axis}
\end{tikzpicture}
\caption{Tre funzioni continue ma non derivabili in $0$: derivate destra e sinistra
diverse (spigolo), rapporto incrementale che tende a $+\infty$ (tangente verticale),
rapporto incrementale che tende a $+\infty$ da destra e a $-\infty$ da sinistra
(cuspide).}
\label{fig:non-derivabili}
\end{figure}

\subsection{Derivabilità e continuità}

Confrontiamo la nozione di derivabilità con quella di continuità. Ricordiamo che una
funzione $f$ è continua in un punto $x$ se (riprendiamo la \eqref{eq:continua-h})
\begin{equation}\label{eq:continua-h-bis}
  \lim_{h\to0}f(x+h)=f(x).
\end{equation}
Gli esempi precedenti mostrano che una funzione continua può non essere derivabile.
Invece vale il viceversa.

\begin{teorema}[La derivabilità implica la continuità]\label{teo:derivabile-continua}
Ogni funzione derivabile in un punto $x$ è continua in $x$.
\end{teorema}

\begin{schema}
\Ipotesi{$f$ derivabile in $x$, cioè esiste finito $f'(x)$.}
\Tesi{$f(x+h)\to f(x)$ per $h\to0$.}
\end{schema}

\begin{strategia}
L'incremento $f(x+h)-f(x)$ si scrive come prodotto del rapporto incrementale, che tende
al numero finito $f'(x)$, per l'incremento $h$, che tende a zero: il prodotto tende a
zero. Il punto essenziale è la \emph{finitezza} della derivata: se il rapporto
incrementale tendesse all'infinito, avremmo una forma indeterminata $\infty\cdot0$.
\end{strategia}

\begin{dimostrazione}
Per $h\neq0$ si ha l'identità
$f(x+h)=f(x)+\frac{f(x+h)-f(x)}{h}\cdot h$; quindi, per le operazioni con i limiti
(\cref{teo:operazioni-funzioni}),
\begin{equation}\label{eq:derivabile-continua}
\begin{split}
  \lim_{h\to0}f(x+h)&=f(x)+\lim_{h\to0}\bigl[f(x+h)-f(x)\bigr]\\
  &=f(x)+\lim_{h\to0}\frac{f(x+h)-f(x)}{h}\cdot\lim_{h\to0}h\\
  &=f(x)+f'(x)\cdot0=f(x). \qedhere
\end{split}
\end{equation}
\end{dimostrazione}

Il teorema non si inverte, come mostrano le funzioni della \cref{fig:non-derivabili};
anzi, la funzione $\sqrt[3]x$ mostra che un rapporto incrementale infinito è compatibile
con la continuità, e la funzione $\sqrt{\abs x}$ che lo sono anche due rapporti
incrementali unilaterali infiniti di segno opposto. In forma contronominale il teorema
fornisce un criterio utile: \emph{una funzione discontinua in un punto non è derivabile
in quel punto}.

\begin{notastorica}[Una funzione continua senza derivata in alcun punto]
Fino a metà Ottocento molti matematici ritenevano che una funzione continua dovesse
essere derivabile tranne al più in punti isolati, come lo spigolo di $\abs x$.
Nel 1872 Karl Weierstrass presentò all'Accademia delle Scienze di Berlino la funzione
\[
  W(x)=\sum_{n=0}^{\infty}a^{n}\cos(b^{n}\pi x), \qquad 0<a<1,\ b \text{ intero dispari},\ ab>1+\tfrac32\pi,
\]
continua su tutto $\R$ ma \emph{priva di derivata in ogni punto}: il suo grafico ha
spigoli a ogni scala. L'esempio, che Bolzano aveva in sostanza anticipato in lavori
rimasti inediti, sconcertò i contemporanei (Hermite parlò di «piaga lamentevole delle
funzioni continue senza derivata»), ma mostrò definitivamente che continuità e
derivabilità sono concetti ben distinti. Un secolo dopo grafici di questo tipo sarebbero
diventati familiari come \emph{frattali}.
\end{notastorica}

\subsection{Il calcolo numerico della derivata}

La definizione suggerisce un modo immediato per calcolare una derivata al computer:
scegliere un $h$ «piccolo» e calcolare il rapporto incrementale. Quanto piccolo? Il
programma che segue lo fa per $f(x)=\sen x$ in $x=1$, dove la derivata esatta vale
$\cos1$ (lo dimostreremo nel \cref{sec:derivate-elementari}), confrontando la
\emph{differenza in avanti} $\frac{f(x+h)-f(x)}{h}$ con la \emph{differenza centrata}
$\frac{f(x+h)-f(x-h)}{2h}$, che ha anch'essa limite $f'(x)$ (\cref{es:differenza-centrata}).

\begin{python}[Derivata numerica di $\sen x$ in $x=1$]
import math

x, esatta = 1.0, math.cos(1.0)        # f = sen, f'(1) = cos 1
print("   h       errore in avanti   errore centrato")
for k in range(1, 13, 1):
    h = 10.0**(-k)
    avanti = (math.sin(x + h) - math.sin(x)) / h
    centrata = (math.sin(x + h) - math.sin(x - h)) / (2 * h)
    ea, ec = abs(avanti - esatta), abs(centrata - esatta)
    print(f"1e-{k:<2d}   {ea:14.3e}   {ec:14.3e}")
\end{python}
\begin{risultato}
   h       errore in avanti   errore centrato
1e-1         4.294e-02        9.001e-04
1e-2         4.216e-03        9.005e-06
1e-3         4.208e-04        9.005e-08
1e-4         4.207e-05        9.004e-10
1e-5         4.207e-06        1.114e-11
1e-6         4.207e-07        2.772e-11
1e-7         4.183e-08        1.943e-10
1e-8         2.970e-09        2.581e-09
1e-9         5.254e-08        2.970e-09
1e-10        5.848e-08        5.848e-08
1e-11        1.169e-06        1.169e-06
1e-12        4.324e-05        1.227e-05
\end{risultato}

\begin{idea}[Troncamento contro arrotondamento: le differenze finite]
L'errore della derivata numerica ha due componenti con andamento opposto: al diminuire
di $h$ la prima decresce, la seconda cresce (\cref{fig:derivata-numerica}).
L'\emph{errore di troncamento} è dovuto al fatto che $h$ non è nullo: per la differenza in avanti è proporzionale ad $h$ (lo si vede
nell'output: dividendo $h$ per $10$ l'errore si divide per $10$), per la differenza
centrata è proporzionale ad $h^{2}$ (l'errore si divide per $100$), come si dimostra con
la formula di Taylor (\cref{es:errore-differenze}). L'\emph{errore di arrotondamento} è dovuto al fatto che il
calcolatore rappresenta $f(x+h)$ e $f(x)$ con circa $16$ cifre significative: la loro
differenza perde tante più cifre quanto più i due numeri sono vicini (è la
cancellazione del \cref{cap:funzioni-continue}), e l'errore assoluto, dell'ordine di
$10^{-16}$, viene poi diviso per $h$. Il risultato è che esiste un passo ottimale, circa
$10^{-8}$ per la differenza in avanti e circa $10^{-5}$ per quella centrata, sotto il
quale \emph{diminuire $h$ peggiora il risultato}. È esattamente il dilemma di chi scrive
un codice alle differenze finite per la fluidodinamica computazionale: raffinare la
griglia riduce l'errore di discretizzazione, ma non all'infinito.
\end{idea}

\begin{figure}[htb]
\centering
\includegraphics[width=0.92\textwidth]{derivata_numerica}
\caption{Errore della derivata numerica di $\sen x$ in $x=1$ in funzione del passo $h$,
in scala bilogaritmica: a destra domina l'errore di troncamento (rette di pendenza $1$ e
$2$), a sinistra quello di arrotondamento, che cresce come $\varepsilon/h$ con
$\varepsilon\approx10^{-16}$ (\texttt{python/fig\_derivata\_numerica.py}).}
\label{fig:derivata-numerica}
\end{figure}

\subsection{Derivate successive}

Se una funzione è derivabile in tutti i punti di un intervallo $(a,b)$, allora la sua
derivata $f'(x)$ è una funzione definita su $(a,b)$. Se questa funzione è a sua volta
derivabile, diremo che la sua derivata $(f')'$ è la \emph{derivata seconda} della
funzione $f$, e indicheremo tale derivata con uno dei simboli
\begin{equation}\label{eq:notazioni-derivata-seconda}
  f'', \qquad \frac{d^{2}f}{dx^{2}}, \qquad D^{2}f, \qquad y'', \qquad \frac{d^{2}y}{dx^{2}},
  \qquad D^{2}y .
\end{equation}
Se a sua volta la derivata seconda è derivabile parleremo di derivata terza $f'''$, e
così via. Useremo il simbolo $f^{(n)}$ per la derivata $n$-esima (con la convenzione
$f^{(0)}=f$). Può accadere che una funzione ammetta derivate fino a un ordine
$n\in\N$; oppure che sia derivabile infinite volte, cioè che ammetta derivate di
qualsiasi ordine. Se la derivata $n$-esima esiste ed è continua in un intervallo, si
dice che $f$ è \emph{di classe $C^{n}$} in quell'intervallo; se $f$ ha derivate di
ogni ordine si dice che è \emph{di classe $C^{\infty}$}. Le funzioni continue si dicono
anche di classe $C^{0}$.

\begin{esempio}\label{es:x-modulo-x}
La funzione $f(x)=x^{2}$ ammette derivate di ogni ordine; infatti abbiamo verificato
che $f'(x)=2x$; poi dalla \eqref{eq:derivata-retta} segue che $f''(x)=2$ e dalla
\eqref{eq:derivata-costante} che $f'''(x)=0$; dato che $f'''$ è costante, si ottiene
$f^{(4)}=0$ e, analogamente, $f^{(n)}(x)=0$ per ogni $n\ge3$.

Invece la funzione $f(x)=x\abs x$ ammette per $x=0$ derivata prima, ma non derivata
seconda. Infatti, essendo
\begin{equation}\label{eq:x-modulo-x}
  f(x)=
  \begin{cases}
    x^{2} & \text{se } x\ge0,\\
    -x^{2} & \text{se } x<0,
  \end{cases}
\end{equation}
$f$ coincide localmente con $x^{2}$ o con $-x^{2}$ in un intorno di ogni $x\neq0$, ed è
quindi derivabile per $x\neq0$, con
\begin{equation}\label{eq:x-modulo-x-derivata}
  f'(x)=
  \begin{cases}
    2x & \text{se } x>0,\\
    -2x & \text{se } x<0,
  \end{cases}
\end{equation}
mentre, se $x=0$, si ha
\begin{equation}\label{eq:x-modulo-x-zero}
  \lim_{h\to0}\frac{f(0+h)-f(0)}{h}=\lim_{h\to0}\frac{h\abs h-0}{h}=\lim_{h\to0}\abs h=0;
\end{equation}
pertanto $f'(0)=0$. In definitiva $f$ è derivabile per ogni $x\in\R$ e la derivata vale
\begin{equation}\label{eq:x-modulo-x-finale}
  f'(x)=2\abs x \qquad\perogni x\in\R .
\end{equation}
Però $f'(x)=2\abs x$ non è derivabile per $x=0$ (\cref{es:modulo-non-derivabile});
quindi non esiste la derivata seconda di $f$ nel punto $x=0$. La funzione $x\abs x$ è di
classe $C^{1}$, ma non è derivabile due volte su $\R$.
\end{esempio}

\needspace{10\baselineskip}
\begin{notastorica}[Da Lagrange a Cauchy: il nome e la definizione]
Né Newton né Leibniz disponevano di una nozione chiara di limite, e il loro modo di
procedere, dividere per un incremento e poi porlo uguale a zero, fu duramente criticato
(si veda la nota storica del \S\ref{sec:def-limite}). Joseph-Louis Lagrange (1736--1813)
tentò di aggirare la difficoltà con l'algebra. Nella \emph{Théorie des fonctions
analytiques} (1797) supponeva che ogni funzione si potesse sviluppare in serie di potenze
dell'incremento, $f(x+i)=f(x)+p\,i+q\,i^{2}+r\,i^{3}+\cdots$, e chiamava il coefficiente
$p$, che è a sua volta una funzione di $x$, \emph{fonction dérivée} di $f$, scrivendolo
$f'x$; ripetendo l'operazione otteneva $f''x$, $f'''x$, \ldots. Da qui vengono il nome
«derivata» e gli apici delle derivate successive. Ma l'esistenza dello sviluppo non si
può giustificare senza sapere già che cosa sia una derivata: il programma conteneva un
circolo vizioso, e ne riparleremo nel \cref{cap:taylor}. La via giusta era quella
indicata da d'Alembert nell'\emph{Encyclopédie}: fondare il calcolo sul limite.
Augustin-Louis Cauchy la percorse fino in fondo nel \emph{Résumé des leçons sur le calcul
infinitésimal} (1823), le lezioni dell'École Polytechnique: la derivata vi è definita, come
nella \cref{def:derivata}, come limite del rapporto incrementale, e di Lagrange restano il
nome e il simbolo. Si capovolge anche l'ordine tra derivata e differenziale: Leibniz
definiva il differenziale per mezzo della tangente e lo trattava di fatto come un
incremento infinitesimo; per Cauchy l'oggetto primitivo è la derivata, $dx$ è una
quantità finita arbitraria e il differenziale è $dy=f'(x)\,dx$, come nel
\cref{sec:retta-tangente}.
\end{notastorica}

%% ----------------------------------------------------------------------------
%%  §5.3  Operazioni con le derivate
%% ----------------------------------------------------------------------------
\section{Operazioni con le derivate}
\label{sec:operazioni-derivate}

Calcolare ogni derivata a partire dalla definizione sarebbe faticoso. Per le derivate
valgono invece regole di calcolo che permettono di derivare funzioni complicate a partire
dalle derivate di poche funzioni elementari.

\begin{teorema}[Operazioni con le derivate]\label{teo:operazioni-derivate}\index{derivata!regole di calcolo}
Se $f$ e $g$ sono due funzioni derivabili in un punto $x$, allora sono derivabili in $x$
anche la somma, la differenza, il prodotto, il quoziente (purché il denominatore sia
diverso da zero), e si ha:
\begin{align}
  &(f\pm g)'=f'\pm g'; \label{eq:derivata-somma}\\
  &(fg)'=f'g+fg'; \label{eq:derivata-prodotto}\\
  &\Bigl(\frac fg\Bigr)'=\frac{f'g-fg'}{g^{2}} \qquad(\text{se } g(x)\neq0). \label{eq:derivata-quoziente}
\end{align}
\end{teorema}

\begin{schema}
\Ipotesi{$f$, $g$ derivabili in $x$; \ per il quoziente, $g(x)\neq0$.}
\Tesi{$f\pm g$, $fg$, $f/g$ derivabili in $x$, con le derivate \eqref{eq:derivata-somma}--\eqref{eq:derivata-quoziente}.}
\end{schema}

\begin{strategia}
In ciascun caso si scrive il rapporto incrementale della funzione composta (somma,
prodotto, quoziente) come combinazione dei rapporti incrementali di $f$ e di $g$, e si
passa al limite con le regole del \cref{teo:operazioni-funzioni}. Per il prodotto e il
quoziente serve un artificio: aggiungere e togliere uno stesso termine, in modo da far
comparire i due rapporti incrementali separatamente. Compariranno anche valori del tipo
$g(x+h)$, che tendono a $g(x)$ perché una funzione derivabile è continua
(\cref{teo:derivabile-continua}): è qui che il teorema del paragrafo precedente si rivela
indispensabile.
\end{strategia}

\begin{dimostrazione}
\emph{Somma.} Dimostriamo la \eqref{eq:derivata-somma} con il segno $+$: per ogni
$h\neq0$ scriviamo il rapporto incrementale relativo alla funzione somma $f+g$,
\begin{equation}\label{eq:rapporto-somma}
\begin{split}
  &\frac{[f(x+h)+g(x+h)]-[f(x)+g(x)]}{h}\\
  &\qquad=\frac{f(x+h)-f(x)}{h}+\frac{g(x+h)-g(x)}{h}.
\end{split}
\end{equation}
Dato che il limite di una somma è uguale alla somma dei limiti, per $h\to0$ si ottiene la
\eqref{eq:derivata-somma}. Il caso con il segno $-$ è identico.

\emph{Prodotto.} Scriviamo il rapporto incrementale relativo alla funzione prodotto $fg$,
aggiungendo e togliendo al numeratore il termine $f(x)g(x+h)$:
\begin{equation}\label{eq:rapporto-prodotto}
\begin{split}
  &\frac{f(x+h)g(x+h)-f(x)g(x)}{h}\\
  &\quad=\frac{f(x+h)g(x+h)-f(x)g(x+h)+f(x)g(x+h)-f(x)g(x)}{h}\\
  &\quad=\frac{f(x+h)-f(x)}{h}\,g(x+h)+f(x)\,\frac{g(x+h)-g(x)}{h}.
\end{split}
\end{equation}
La funzione $g$, essendo per ipotesi derivabile in $x$, è anche continua
(\cref{teo:derivabile-continua}); quindi al limite per $h\to0$ risulta $g(x+h)\to g(x)$.
Dalla relazione sopra scritta si ottiene la tesi passando al limite per $h\to0$.

\emph{Quoziente.} Supponiamo $g(x)\neq0$. Poiché $g$ è continua in $x$, per il teorema
della permanenza del segno (\cref{teo:permanenza-segno-continue}, applicato a $g$ oppure
a $-g$) esiste un numero $\delta>0$ per cui, se $\abs h<\delta$, allora $g(x+h)\neq0$:
la funzione $f/g$ è dunque definita in un intorno di $x$. Scriviamo il rapporto
incrementale di $f/g$, aggiungendo e togliendo al numeratore $f(x)g(x)$:
\begin{equation}\label{eq:rapporto-quoziente}
\begin{split}
  &\Bigl(\frac{f(x+h)}{g(x+h)}-\frac{f(x)}{g(x)}\Bigr)\frac1h
  =\frac{f(x+h)g(x)-f(x)g(x+h)}{g(x+h)\,g(x)\,h}\\
  &\quad=\frac{f(x+h)g(x)-f(x)g(x)+f(x)g(x)-f(x)g(x+h)}{g(x+h)\,g(x)\,h}\\
  &\quad=\frac1{g(x+h)\,g(x)}\Bigl(\frac{f(x+h)-f(x)}{h}\,g(x)\\
  &\qquad\qquad-f(x)\,\frac{g(x+h)-g(x)}{h}\Bigr).
\end{split}
\end{equation}
Al limite per $h\to0$ si ottiene la tesi, ricordando che, come nel caso del prodotto, la
funzione $g$ è continua in $x$ e quindi $g(x+h)$ tende a $g(x)\neq0$.
\end{dimostrazione}

Notiamo che un caso particolarmente importante di derivazione di un prodotto si ha
quando una delle due funzioni è costante. Dato che la derivata di una costante vale zero,
dalla regola \eqref{eq:derivata-prodotto} si ottiene
\begin{equation}\label{eq:derivata-costante-per}
  (cf)'=cf' \qquad(c \text{ costante}).
\end{equation}
Insieme alla \eqref{eq:derivata-somma}, ciò significa che la derivazione è
un'operazione \emph{lineare}: $(\alpha f+\beta g)'=\alpha f'+\beta g'$ per ogni coppia di
costanti $\alpha$, $\beta$. Un altro caso particolare utile è quello del reciproco, che si
ottiene dalla \eqref{eq:derivata-quoziente} con $f=1$:
\begin{equation}\label{eq:derivata-reciproco}
  \Bigl(\frac1g\Bigr)'=-\frac{g'}{g^{2}} \qquad(\text{se } g(x)\neq0).
\end{equation}

\begin{osservazione}[Prodotto di più fattori]\label{oss:prodotto-n}
Applicando ripetutamente la \eqref{eq:derivata-prodotto} si ottiene, per tre fattori,
$(fgh)'=f'gh+fg'h+fgh'$ e, per induzione, la regola generale: la derivata di un prodotto
di $n$ funzioni è la somma degli $n$ prodotti che si ottengono derivando, a turno, un
solo fattore e lasciando invariati gli altri.
\end{osservazione}

\begin{attenzione}[La derivata di un prodotto non è il prodotto delle derivate]
L'errore più frequente è scrivere $(fg)'=f'g'$. Basta l'esempio $f(x)=g(x)=x$: si ha
$fg=x^{2}$, la cui derivata è $2x$, mentre $f'g'=1\cdot1=1$. La regola corretta,
\eqref{eq:derivata-prodotto}, fu trovata da Leibniz; nei suoi manoscritti del 1675 egli
si chiese dapprima se il differenziale di un prodotto fosse il prodotto dei
differenziali, e lo escluse proprio verificandolo su un esempio.
\end{attenzione}

\begin{idea}[Ritorno all'efficienza massima]
Nel \cref{cap:funzioni-continue} abbiamo trovato il massimo dell'efficienza
$E(C_{L})=\frac{C_{L}}{C_{D0}+kC_{L}^{2}}$ di una polare parabolica con una disuguaglianza
algebrica. Con la regola del quoziente calcoliamo la derivata:
\[
  E'(C_{L})=\frac{1\cdot(C_{D0}+kC_{L}^{2})-C_{L}\cdot2kC_{L}}{(C_{D0}+kC_{L}^{2})^{2}}
  =\frac{C_{D0}-kC_{L}^{2}}{(C_{D0}+kC_{L}^{2})^{2}} .
\]
Il denominatore è positivo; il numeratore è positivo per $C_{L}<\sqrt{C_{D0}/k}$, nullo
per $C_{L}=C_{L}^{*}=\sqrt{C_{D0}/k}$, negativo per $C_{L}>C_{L}^{*}$. La derivata, cioè la
pendenza del grafico della \cref{fig:efficienza}, si annulla esattamente nel punto di
massimo trovato allora, ed è positiva prima e negativa dopo. Non è una coincidenza: nel
prossimo capitolo dimostreremo (teorema di Fermat) che nei punti di massimo o di minimo
interni la derivata di una funzione derivabile si annulla, e (teorema di Lagrange) che il
segno della derivata determina la crescenza o la decrescenza. Il calcolo differenziale
trasformerà la ricerca dei massimi, che il teorema di Weierstrass sa solo garantire,
in un procedimento sistematico.
\end{idea}

\begin{notastorica}[Le regole della \emph{Nova methodus}]
Il titolo della memoria di Leibniz del 1684, nell'epigrafe del capitolo, promette un
metodo che «non si arresta davanti alle quantità fratte o irrazionali». I procedimenti
precedenti per le tangenti e per i massimi, come quello di Fermat (si veda la nota storica
del \S\ref{sec:fermat}), andavano applicati all'intera espressione della funzione, dalla
quale occorreva spesso eliminare prima frazioni e radicali. La novità fu riconoscere
nella differenziazione un'\emph{operazione} con regole proprie: il differenziale di una
somma, di un prodotto, di un quoziente, di una potenza si ottiene da quelli delle parti, e
un'espressione complicata si tratta un pezzo alla volta. La memoria enunciava le regole
senza dimostrarle; altrove Leibniz giustificava quella del prodotto scrivendo
$d(xy)=(x+dx)(y+dy)-xy=x\,dy+y\,dx+dx\,dy$ e trascurando $dx\,dy$, infinitamente
piccolo rispetto agli altri termini. Nella
\eqref{eq:rapporto-prodotto} quel termine non si trascura: scrivendo
$g(x+h)=g(x)+[g(x+h)-g(x)]$, esso compare come
$\frac{f(x+h)-f(x)}{h}\,[g(x+h)-g(x)]$, che tende a $f'(x)\cdot0=0$ perché $g$ è continua.
È il limite a giustificare ciò che Leibniz faceva per intuizione. Per le radici bastava
un accorgimento: se $z=x^{1/k}$, allora $x=z^{k}$, $dx=kz^{k-1}\,dz$ e quindi
$dz=\frac1k\,x^{1/k-1}\,dx$, che è in sostanza la regola della funzione inversa
(\cref{teo:derivata-inversa}). Tra gli esempi della memoria c'è la legge della rifrazione
della luce, ricavata come problema di minimo (\cref{es:snell}).
\end{notastorica}

%% ----------------------------------------------------------------------------
%%  §5.4  Derivate delle funzioni composte e delle funzioni inverse
%% ----------------------------------------------------------------------------
\section{Derivate delle funzioni composte e delle funzioni inverse}
\label{sec:composte-inverse}

Una delle più importanti regole di derivazione è quella relativa alle funzioni composte.
Se $y$ è funzione di $z$, $y=f(z)$, e $z$ è a sua volta funzione di $x$, $z=g(x)$, allora
$y=f(g(x))$ è la \emph{funzione composta} risultante; si usa anche il simbolo
$f(g(x))=(f\circ g)(x)$. Sono funzioni composte, ad esempio, $y=\sen x^{2}$ (con
$y=\sen z$, $z=x^{2}$), oppure $y=\sen^{2}x$ (con $y=z^{2}$, $z=\sen x$).

\begin{teorema}[Derivazione delle funzioni composte]\label{teo:derivata-composta}\index{derivata!di una funzione composta}\index{regola della catena}
Sia $g$ una funzione definita in un intorno di $x$ e derivabile in $x$, e sia $f$ una
funzione definita in un intorno del punto $g(x)$ e derivabile in $g(x)$. Allora la
funzione composta $f(g(x))$ è derivabile in $x$, e si ha
\begin{equation}\label{eq:derivata-composta}
  D f(g(x))=f'(g(x))\cdot g'(x).
\end{equation}
\end{teorema}

\begin{schema}
\Ipotesi{$g$ derivabile in $x$; \ $f$ derivabile in $y=g(x)$.}
\Tesi{$(f\circ g)'(x)=f'(g(x))\,g'(x)$.}
\end{schema}

\begin{strategia}
L'idea naturale è moltiplicare e dividere il rapporto incrementale della composta per
l'incremento $k=g(x+h)-g(x)$ della funzione interna:
\[
  \frac{f(g(x+h))-f(g(x))}{h}=\frac{f(g(x+h))-f(g(x))}{g(x+h)-g(x)}\cdot\frac{g(x+h)-g(x)}{h}.
\]
Il secondo fattore tende a $g'(x)$, il primo è un rapporto incrementale di $f$ nel punto
$g(x)$, con incremento $k\to0$. Il procedimento è corretto se $k\neq0$ per $h\neq0$
piccolo; ma può accadere che $g(x+h)=g(x)$ per valori di $h$ arbitrariamente vicini a
zero (si pensi a una $g$ costante, o a $g(x)=x^{2}\sen\frac1x$ in $x=0$, che si annulla in
tutti i punti $\frac1{n\pi}$), e allora la divisione per $k$ non è lecita. Per superare la
difficoltà si sostituisce il rapporto incrementale di $f$ con una funzione $F(k)$ che
coincide con esso per $k\neq0$ e che, per $k=0$, vale proprio $f'(g(x))$: così definita,
$F$ è continua in $0$, e la continuità permette di passare al limite anche lungo valori
$k=0$ (\cref{cor:limite-composta-continua}).
\end{strategia}

\begin{dimostrazione}
\emph{Un caso particolare.} Per semplificare, consideriamo preliminarmente il caso in
cui risulti $g(x+h)\neq g(x)$ per ogni $h\neq0$ abbastanza piccolo. Il rapporto
incrementale della funzione composta, nel punto $x$, vale
\begin{equation}\label{eq:composta-caso-particolare}
  \frac{f(g(x+h))-f(g(x))}{h}=\frac{f(g(x+h))-f(g(x))}{g(x+h)-g(x)}\cdot\frac{g(x+h)-g(x)}{h}.
\end{equation}
Nel primo dei due quozienti a secondo membro compare il rapporto incrementale della
funzione $f$ nel punto $g(x)$, con incremento $k=g(x+h)-g(x)$. Tale incremento $k$ tende a
zero per $h\to0$, dato che $g$ è continua in $x$ (\cref{teo:derivabile-continua}), ed è
diverso da zero per $h\neq0$; quindi, per il teorema sul limite delle funzioni composte
(\cref{teo:limite-composta}),
\begin{equation}\label{eq:composta-limite-particolare}
\begin{split}
  \lim_{h\to0}\frac{f(g(x+h))-f(g(x))}{g(x+h)-g(x)}
  &=\lim_{k\to0}\frac{f(g(x)+k)-f(g(x))}{k}\\
  &=f'(g(x)),
\end{split}
\end{equation}
e la \eqref{eq:composta-caso-particolare} dà la tesi \eqref{eq:derivata-composta}.

\emph{Il caso generale.} Osserviamo anzitutto che, essendo $g$ continua in $x$, esiste
$\delta>0$ tale che $g(x+h)$ appartiene all'intorno di $g(x)$ in cui $f$ è definita per
ogni $\abs h<\delta$: la funzione composta è dunque definita in un intorno di $x$, e
tutti gli incrementi $k$ che considereremo stanno nel dominio della funzione $F$ che
stiamo per introdurre. Poniamo $y=g(x)$ e
\begin{equation}\label{eq:funzione-F}
  F(k)=
  \begin{cases}
    \dfrac{f(y+k)-f(y)}{k} & \text{se } k\neq0,\\[1.5ex]
    f'(y) & \text{se } k=0,
  \end{cases}
\end{equation}
definita per $k$ in un intorno di $0$ (quello per cui $y+k$ sta nel dominio di $f$). Per
l'ipotesi di derivabilità di $f$ nel punto $y$ e per la definizione di $F(0)$ risulta
\begin{equation}\label{eq:F-continua}
  \lim_{k\to0}F(k)=f'(y)=F(0),
\end{equation}
cioè $F$ è continua nel punto $k=0$. Posto
\begin{equation}\label{eq:k-di-h}
  k=k(h)=g(x+h)-g(x),
\end{equation}
ed essendo $g(x)=y$, $g(x+h)=y+k$, per ogni $h\neq0$ con $k\neq0$ risulta
\begin{equation}\label{eq:composta-identita}
\begin{split}
  \frac{f(g(x+h))-f(g(x))}{h}&=\frac{f(y+k)-f(y)}{k}\cdot\frac kh\\
  &=F(k)\cdot\frac{g(x+h)-g(x)}{h}.
\end{split}
\end{equation}
La novità rispetto alla dimostrazione precedente è che l'identità
\begin{equation}\label{eq:composta-identita-F}
  \frac{f(g(x+h))-f(g(x))}{h}=F(k(h))\cdot\frac{g(x+h)-g(x)}{h}
\end{equation}
vale non solo per $k\neq0$, ma anche per $k=0$, perché in tal caso, essendo
$k=g(x+h)-g(x)=0$, risulta nullo il secondo membro della \eqref{eq:composta-identita-F},
ma anche il primo membro, dato che $g(x+h)=g(x)$. Passiamo al limite nella
\eqref{eq:composta-identita-F} per $h\to0$. Per la continuità di $g$ in $x$, l'incremento
$k(h)$ converge a zero per $h\to0$; allora, per la continuità di $F$ in $0$
\eqref{eq:F-continua} e per il \cref{cor:limite-composta-continua}, $F(k(h))\to F(0)$, e
si ottiene
\begin{equation}\label{eq:composta-conclusione}
\begin{split}
  \lim_{h\to0}\frac{f(g(x+h))-f(g(x))}{h}&=\lim_{h\to0}F(k(h))\cdot\lim_{h\to0}\frac{g(x+h)-g(x)}{h}\\
  &=F(0)\,g'(x)=f'(g(x))\,g'(x). \qedhere
\end{split}
\end{equation}
\end{dimostrazione}

Così, ad esempio, in base alla regola di derivazione delle funzioni composte, la derivata
della funzione $y=\sen x^{2}$ vale $y'=\cos(x^{2})\cdot2x$, mentre la derivata della
funzione $y=\sen^{2}x=(\sen x)^{2}$ vale $y'=2\sen x\cdot\cos x$ (usiamo la derivata di
$\sen$, che calcoleremo nel \cref{sec:derivate-elementari}, e quella di $z^{2}$, già nota
dalla \eqref{eq:tasso-t2}). Con la notazione di Leibniz, posto $y=f(z)$ e $z=g(x)$, la
\eqref{eq:derivata-composta} assume la forma suggestiva
\[
  \frac{dy}{dx}=\frac{dy}{dz}\cdot\frac{dz}{dx},
\]
come se i differenziali $dz$ si «semplificassero»: per questo si parla anche di
\emph{regola della catena}. Si tratta però di una regola mnemonica, non di una
dimostrazione, che ha dovuto superare proprio la difficoltà di un incremento $k$ della
variabile intermedia $z$ che può annullarsi.

\begin{osservazione}\label{oss:caratheodory}
L'artificio della funzione $F$ ha un significato che va oltre questa dimostrazione. La
\eqref{eq:funzione-F} dice che \emph{$f$ è derivabile in $y$ se e solo se esiste una
funzione $F$, continua in $0$, tale che $f(y+k)-f(y)=F(k)\,k$ per ogni $k$ in un intorno
di $0$}; e allora
$f'(y)=F(0)$. In questa forma, dovuta a Carathéodory, la derivabilità appare come la
possibilità di «fattorizzare» l'incremento della funzione in un fattore continuo per
l'incremento della variabile, e la regola della catena si riduce al fatto che il prodotto
e la composizione di funzioni continue sono continui.
\end{osservazione}

\begin{idea}[La pressione dinamica durante un'accelerazione]
Un velivolo accelera in volo orizzontale, a quota costante (densità $\rho$ costante),
con velocità $V(t)$. La pressione dinamica $q=\frac12\rho V^{2}$ è funzione composta del
tempo, $q(t)=f(V(t))$ con $f(V)=\frac12\rho V^{2}$, e la regola della catena dà
\[
  \frac{dq}{dt}=\frac{dq}{dV}\cdot\frac{dV}{dt}=\rho V\,\frac{dV}{dt}.
\]
Il tasso di crescita della pressione dinamica, e quindi dei carichi aerodinamici a
parità di coefficienti, non dipende solo dall'accelerazione $\frac{dV}{dt}$, ma anche
dalla velocità: la stessa accelerazione di $1$~m/s$^{2}$ fa crescere $q$, al livello del
mare ($\rho=1{,}225$~kg/m$^{3}$), di circa $61$~Pa al secondo a $50$~m/s e di circa
$245$~Pa al secondo a $200$~m/s.
\end{idea}

\subsection{La derivata della funzione inversa}

Esaminiamo ora la regola di derivazione delle funzioni inverse. Ricordiamo quanto già
detto circa le funzioni strettamente monotòne: una funzione $f$ è \emph{strettamente
crescente} nell'intervallo $[a,b]$ se
\begin{equation}\label{eq:strettamente-crescente}
  a\le x_{1}<x_{2}\le b \quad\implica\quad f(x_{1})<f(x_{2}).
\end{equation}
Se $f$ è continua e strettamente crescente in $[a,b]$ allora è anche \emph{invertibile},
cioè a ogni $y\in[f(a),f(b)]$ corrisponde un solo $x\in[a,b]$ per cui $f(x)=y$, e si indica
$x=f^{-1}(y)$; la stessa proprietà vale per le funzioni strettamente decrescenti
(\cref{teo:invertibilita}). Inoltre la funzione inversa è continua
(\cref{teo:continuita-inversa}).

\needspace{9\baselineskip}
\begin{teorema}[Derivazione delle funzioni inverse]\label{teo:derivata-inversa}\index{derivata!della funzione inversa}
Sia $f$ una funzione continua e strettamente crescente (oppure strettamente decrescente)
in un intervallo $[a,b]$. Se $f$ è derivabile in un punto $x\in(a,b)$ e se $f'(x)\neq0$,
allora anche $f^{-1}$ è derivabile nel punto $y=f(x)$ e la derivata vale
\begin{equation}\label{eq:derivata-inversa}
  D f^{-1}(y)=\frac1{f'(x)}=\frac1{f'(f^{-1}(y))}.
\end{equation}
\end{teorema}

\begin{schema}
\Ipotesi{$f$ continua e strettamente monotòna in $[a,b]$; \ $x\in(a,b)$; \ $f$ derivabile
in $x$; \ $f'(x)\neq0$.}
\Tesi{$f^{-1}$ derivabile in $y=f(x)$, \ $(f^{-1})'(y)=1/f'(x)$.}
\end{schema}

\begin{figure}[htb]
\centering
\begin{tikzpicture}
\begin{axis}[width=0.9\textwidth, height=5.8cm, xmin=0, xmax=6.2, ymin=0, ymax=4.4,
  axis lines=left, xtick={1.8,4.4}, xticklabels={{$x=f^{-1}(y)$},{$x+h=f^{-1}(y+k)$}},
  ytick={1.4,3.25}, yticklabels={{$y=f(x)$},{$y+k=f(x+h)$}},
  tick label style={font=\scriptsize}, clip=false]
  \addplot[curva, domain=0.6:6, samples=100, smooth]
    coordinates {(0.6,0.7) (1.8,1.4) (2.8,2.3) (3.6,2.6) (4.4,3.25) (5.2,3.85) (6,4.1)};
  \addplot[tratteggio] coordinates {(0,1.4) (1.8,1.4) (1.8,0)};
  \addplot[tratteggio] coordinates {(0,3.25) (4.4,3.25) (4.4,0)};
  \draw[decorate, decoration={brace, amplitude=5pt}, grigio] (axis cs:1.8,0.3) -- node[above=5pt, font=\scriptsize, text=black] {$h$} (axis cs:4.4,0.3);
  \draw[decorate, decoration={brace, amplitude=5pt, mirror}, grigio] (axis cs:0.3,1.4) -- node[right=5pt, font=\scriptsize, text=black] {$k$} (axis cs:0.3,3.25);
  \addplot[only marks, mark=*, mark size=1.6pt, rame] coordinates {(1.8,1.4) (4.4,3.25)};
\end{axis}
\end{tikzpicture}
\caption{Il rapporto incrementale della funzione inversa è il reciproco di quello di $f$:
all'incremento $h$ della $x$ corrisponde l'incremento $k$ della $y$, e
$\frac{f^{-1}(y+k)-f^{-1}(y)}{k}=\frac hk=\frac1{k/h}$. Quando $k\to0$ anche $h\to0$,
perché $f^{-1}$ è continua.}
\label{fig:derivata-inversa}
\end{figure}

\begin{strategia}
Il grafico di $f^{-1}$ è quello di $f$ visto «scambiando gli assi», e la pendenza di una
retta, scambiando gli assi, diventa la sua reciproca. In formule: all'incremento $k$ della
variabile $y$ corrisponde l'incremento $h=f^{-1}(y+k)-f^{-1}(y)$ della variabile $x$, e il
rapporto incrementale di $f^{-1}$ è $\frac hk$, reciproco del rapporto incrementale
$\frac kh$ di $f$ (\cref{fig:derivata-inversa}). Per passare al limite servono due fatti:
che $h\neq0$ quando $k\neq0$ (iniettività) e che $h\to0$ quando $k\to0$ (continuità della
funzione inversa, che abbiamo dimostrato nel \cref{sec:monotone-inverse}).
\end{strategia}

\begin{dimostrazione}
Con riferimento alla \cref{fig:derivata-inversa}, a $x$ corrisponde $y=f(x)$, che è
interno all'intervallo immagine di estremi $f(a)$, $f(b)$; per $k\neq0$ abbastanza
piccolo anche $y+k$ sta in tale intervallo, e poniamo
\[
  h=h(k)=f^{-1}(y+k)-f^{-1}(y),
\]
cosicché $x+h=f^{-1}(y+k)$ e $y+k=f(x+h)$, cioè $k=f(x+h)-f(x)$. Il rapporto incrementale
relativo a $f^{-1}$ vale quindi
\begin{equation}\label{eq:rapporto-inversa}
  \frac{f^{-1}(y+k)-f^{-1}(y)}{k}=\frac{h}{f(x+h)-f(x)} .
\end{equation}
Dato che $f$ è strettamente monotòna, risulta $h\neq0$ se e solo se $k\neq0$. Inoltre, se
$k$ tende a zero allora anche $h$ tende a zero; infatti
$h=f^{-1}(y+k)-f^{-1}(y)\to0$ per $k\to0$, perché $f^{-1}$ è una funzione continua
(\cref{teo:continuita-inversa}). Dato che a secondo membro della
\eqref{eq:rapporto-inversa} compare il reciproco del rapporto incrementale di $f$, che
tende a $f'(x)\neq0$ quando l'incremento $h$ tende a zero, passando al limite nella
\eqref{eq:rapporto-inversa} per $k\to0$ (\cref{teo:limite-composta}, con la funzione
interna $h(k)$, che tende a $0$ ed è diversa da $0$ per $k\neq0$) si ottiene
\[
  \lim_{k\to0}\frac{f^{-1}(y+k)-f^{-1}(y)}{k}=\lim_{h\to0}\frac1{\dfrac{f(x+h)-f(x)}{h}}=\frac1{f'(x)}. \qedhere
\]
\end{dimostrazione}

\begin{strategia}[Un'altra via]
È la via seguita da Giusti \cite{giusti}, che ripete l'artificio usato per la funzione
composta. Invece di controllare che l'incremento $h(k)$ non si annulli, come richiede il
\cref{teo:limite-composta}, si assegna in $h=0$ al reciproco del rapporto incrementale di
$f$ il valore $\frac1{f'(x)}$, che lo rende continuo. Il rapporto incrementale di
$f^{-1}$ diventa allora la composizione di due funzioni continue, e per passare al limite
basta il \cref{cor:limite-composta-continua}.
\end{strategia}

\begin{dimostrazione}[Seconda dimostrazione]
Per gli $h$ tali che $x+h\in[a,b]$ poniamo
\begin{equation}\label{eq:funzione-G-inversa}
  G(h)=
  \begin{cases}
    \dfrac{h}{f(x+h)-f(x)} & \text{se } h\neq0,\\[1.5ex]
    \dfrac1{f'(x)} & \text{se } h=0;
  \end{cases}
\end{equation}
per $h\neq0$ il denominatore non si annulla, perché $f$ è strettamente monotòna. Il
rapporto incrementale di $f$ tende a $f'(x)\neq0$, e quindi, per il limite del reciproco
(\cref{teo:operazioni-funzioni}), $G(h)\to\frac1{f'(x)}=G(0)$ per $h\to0$: $G$ è continua
in $0$. Sia ora, come nella prima dimostrazione, $h(k)=f^{-1}(y+k)-f^{-1}(y)$, così che
$f(x+h(k))=y+k$. Per $k\neq0$ si ha $h(k)\neq0$, e dunque
\[
  \frac{f^{-1}(y+k)-f^{-1}(y)}{k}=\frac{h(k)}{f(x+h(k))-f(x)}=G(h(k)).
\]
Per la continuità di $f^{-1}$ (\cref{teo:continuita-inversa}) si ha $h(k)\to0$ per
$k\to0$; poiché $G$ è continua in $0$, il \cref{cor:limite-composta-continua} dà
$G(h(k))\to G(0)=\frac1{f'(x)}$, che è la tesi.
\end{dimostrazione}

\begin{osservazione}[La funzione inversa secondo Carathéodory]
Nel linguaggio dell'\cref{oss:caratheodory} la seconda dimostrazione ha una lettura
particolarmente semplice. L'incremento di $f$ si fattorizza come $f(x+h)-f(x)=F(h)\,h$,
con $F$ continua in $0$ e $F(0)=f'(x)\neq0$; inoltre $F(h)\neq0$ anche per $h\neq0$, perché
$f$ è iniettiva. Sostituendo $h=h(k)$ si ottiene $k=F(h(k))\,h(k)$, cioè
\[
  f^{-1}(y+k)-f^{-1}(y)=h(k)=\frac1{F(h(k))}\cdot k,
\]
e il fattore $\frac1{F(h(k))}=G(h(k))$ è continuo in $k=0$. L'incremento della funzione
inversa ammette dunque una fattorizzazione dello stesso tipo, con il fattore reciproco:
per questo la derivata dell'inversa è il reciproco della derivata di $f$. Come la regola
della catena, anche la regola della funzione inversa si riduce così a proprietà delle
funzioni continue: qui la continuità del reciproco e della composizione.
\end{osservazione}

Applichiamo il teorema a un esempio concreto: la funzione $y=f(x)=x^{2}$ è continua e
strettamente crescente per $x>0$ (su ogni intervallo $[a,b]$ con $0<a<b$). La funzione
inversa è $x=f^{-1}(y)=\sqrt y$. Abbiamo già visto che la funzione $y=x^{2}$ è derivabile
e che la derivata vale $y'=2x\neq0$ per $x>0$. In base al teorema di derivazione delle
funzioni inverse, anche $x=\sqrt y$ è derivabile per $y>0$ e la derivata vale
\begin{equation}\label{eq:derivata-radice}
  D\sqrt y=\frac1{D(x^{2})}=\frac1{2x}=\frac1{2\sqrt y}.
\end{equation}
(Che il teorema, enunciato su un intervallo chiuso, valga in ogni punto $y>0$ segue dal
carattere locale della derivata: basta scegliere $[a,b]$ in modo che $x=\sqrt y$ sia
interno.)

\begin{osservazione}[Quando $f'(x)=0$]\label{oss:inversa-derivata-nulla}
L'ipotesi $f'(x)\neq0$ è necessaria: se $f'(x)=0$, la funzione inversa \emph{non} è
derivabile in $y=f(x)$. Infatti, se lo fosse, applicando il teorema di derivazione delle
funzioni composte all'identità $f^{-1}(f(t))=t$ nel punto $t=x$ si otterrebbe
$(f^{-1})'(y)\cdot f'(x)=1$, cioè $(f^{-1})'(y)\cdot0=1$, che è assurdo. Ad esempio
$f(x)=x^{3}$ è strettamente crescente e derivabile, con $f'(0)=0$; la sua inversa
$\sqrt[3]y$ non è derivabile in $0$, dove il grafico ha tangente verticale
(\cref{es:derivate-infinite}). Geometricamente, la retta tangente orizzontale al grafico
di $f$ diventa, scambiando gli assi, una tangente verticale al grafico di $f^{-1}$.
Con la notazione di Leibniz la \eqref{eq:derivata-inversa} si scrive
$\frac{dx}{dy}=1\Big/\frac{dy}{dx}$.
\end{osservazione}

%% ----------------------------------------------------------------------------
%%  §5.5  Derivate delle funzioni elementari
%% ----------------------------------------------------------------------------
\section{Derivate delle funzioni elementari}
\label{sec:derivate-elementari}

In questo paragrafo calcoliamo le derivate delle funzioni elementari. Insieme alle regole
dei due paragrafi precedenti, esse permettono di derivare tutte le funzioni che si
ottengono da quelle elementari mediante somme, prodotti, quozienti, composizioni e
inversioni.

\subsection{Potenze con esponente naturale e polinomi}

\begin{proposizione}\label{prop:derivata-xn}
Per ogni $n\in\N$ e ogni $x\in\R$ si ha
\begin{equation}\label{eq:derivata-xn}
  Dx^{n}=nx^{n-1}
\end{equation}
(per $n=1$ si intende $Dx=1$).
\end{proposizione}

\begin{schema}
\Ipotesi{$n\in\N$; \ $x\in\R$.}
\Tesi{$x^{n}$ è derivabile in $x$ e $Dx^{n}=nx^{n-1}$.}
\end{schema}

\begin{strategia}
La formula dipende da un indice $n\in\N$: è naturale procedere per induzione. Il caso
$n=1$ è la derivata della funzione lineare $x$, già calcolata nell'\cref{es:derivata-lineare};
nel passo induttivo si scrive $x^{n+1}$ come prodotto $x^{n}\cdot x$ e si applica la regola
del prodotto \eqref{eq:derivata-prodotto}, usando l'ipotesi induttiva per la derivata del
primo fattore.
\end{strategia}

\begin{dimostrazione}
Questa formula (che per $x>0$ si ottiene anche dalla più generale \eqref{eq:derivata-xb}) si può
dimostrare facendo uso del principio di induzione (\cref{teo:induzione}): abbiamo già
verificato che $Dx=1$ \eqref{eq:derivata-retta}; quindi la \eqref{eq:derivata-xn} è vera
per $n=1$. Supponiamo, secondo lo schema del principio di induzione, che la
\eqref{eq:derivata-xn} sia vera per un certo $n$, e calcoliamo per mezzo della regola di
derivazione del prodotto:
\begin{equation}\label{eq:derivata-xn-passo}
\begin{split}
  Dx^{n+1}&=D(x^{n}\cdot x)=D(x^{n})\,x+x^{n}\,Dx\\
  &=nx^{n-1}\,x+x^{n}\cdot1=(n+1)x^{n}.
\end{split}
\end{equation}
Abbiamo quindi verificato che la \eqref{eq:derivata-xn} vale anche per l'indice $n+1$.
Perciò la \eqref{eq:derivata-xn} è provata per ogni $n\in\N$.
\end{dimostrazione}

\begin{strategia}[Un'altra via]
Al posto dell'induzione, una formula chiusa, come in Giusti \cite{giusti}. La
scomposizione \eqref{eq:scomposizione-potenze} di $u^{n}-v^{n}$ permette di semplificare
il fattore $\xi-x$ nel rapporto incrementale \eqref{eq:derivata-xi}, che diventa una
somma di $n$ monomi; ciascuno tende a $x^{n-1}$, e il coefficiente $n$ della formula
conta proprio questi addendi. La regola del prodotto non serve.
\end{strategia}

\begin{dimostrazione}[Seconda dimostrazione]
La \eqref{eq:scomposizione-potenze}, ricavata nel \cref{lem:stima-potenze} per
$0\le v<u$, vale per ogni coppia di numeri reali $u$, $v$ (con la convenzione $t^{0}=1$
per ogni $t$). Detta infatti $S=u^{n-1}+u^{n-2}v+\dots+uv^{n-2}+v^{n-1}$ la somma tra
parentesi, si ha
\[
  uS=u^{n}+u^{n-1}v+\dots+uv^{n-1}, \qquad vS=u^{n-1}v+\dots+uv^{n-1}+v^{n},
\]
e nella differenza $uS-vS$ tutti i termini si cancellano, tranne $u^{n}-v^{n}$. Con
$u=\xi$, $v=x$ e $\xi\neq x$ si ottiene
\[
  \frac{\xi^{n}-x^{n}}{\xi-x}=\xi^{n-1}+\xi^{n-2}x+\dots+\xi x^{n-2}+x^{n-1}
  =\sum_{k=0}^{n-1}\xi^{n-1-k}x^{k}.
\]
Per $\xi\to x$ ciascuno degli $n$ addendi tende a $x^{n-1-k}x^{k}=x^{n-1}$, per il limite
di un prodotto (\cref{teo:operazioni-funzioni}); quindi il rapporto incrementale tende a
$nx^{n-1}$, e per la \eqref{eq:derivata-xi} si ha $Dx^{n}=nx^{n-1}$.
\end{dimostrazione}

Notiamo che il risultato ottenuto, insieme alla linearità della derivata, ci permette di
calcolare la derivata di un \emph{polinomio} qualsiasi:
\begin{equation}\label{eq:derivata-polinomio}
\begin{gathered}
  y=a_{n}x^{n}+a_{n-1}x^{n-1}+\dots+a_{1}x+a_{0};\\
  y'=na_{n}x^{n-1}+(n-1)a_{n-1}x^{n-2}+\dots+a_{1}.
\end{gathered}
\end{equation}
In particolare, la derivata di un polinomio di grado $n\in\N$ è un polinomio di grado
$n-1$, e un polinomio è derivabile infinite volte (è di classe $C^{\infty}$). Per gli
esponenti interi negativi la regola del reciproco \eqref{eq:derivata-reciproco} dà, per
$x\neq0$, $Dx^{-n}=D\frac1{x^{n}}=-\frac{nx^{n-1}}{x^{2n}}=-nx^{-n-1}$: la
\eqref{eq:derivata-xn} vale per ogni esponente intero.

\subsection{Logaritmi ed esponenziali}

\begin{proposizione}[Derivata del logaritmo]\label{prop:derivata-log}
Per ogni base $a>0$, $a\neq1$, e per ogni $x>0$ si ha
\begin{equation}\label{eq:derivata-loga}
  D\log_{a}x=\frac1x\log_{a}e=\frac1{x\log a};
\end{equation}
in particolare, per il logaritmo naturale,
\begin{equation}\label{eq:derivata-log}
  D\log x=\frac1x \qquad\perogni x>0 .
\end{equation}
\end{proposizione}

\begin{schema}
\Ipotesi{$a>0$, \ $a\neq1$; \ $x>0$.}
\Tesi{$\log_{a}$ è derivabile in $x$ e $D\log_{a}x=\frac1{x\log a}$; \ in particolare
$D\log x=\frac1x$.}
\end{schema}

\begin{strategia}
Le proprietà del logaritmo trasformano la differenza $\log(x+h)-\log x$ nel logaritmo di
un rapporto, $\log\bigl(1+\frac hx\bigr)$; il rapporto incrementale diventa allora,
a meno del fattore $\frac1x$, la quantità $\frac{\log(1+t)}{t}$ con $t=\frac hx\to0$, di cui
conosciamo il limite dal \cref{cap:funzioni-continue}. È qui che si vede perché i
matematici preferiscono la base $e$: è l'unica per cui quel limite vale $1$.
\end{strategia}

\begin{dimostrazione}
Sia $x>0$ e sia $0<\abs h<x$, così che $x+h>0$. Per la \eqref{eq:log2}
\begin{equation}\label{eq:rapporto-log}
  \frac{\log(x+h)-\log x}{h}=\frac1h\log\frac{x+h}{x}
  =\frac1x\cdot\frac{\log\bigl(1+\frac hx\bigr)}{\frac hx}.
\end{equation}
La funzione $t=g(h)=\frac hx$ tende a $0$ per $h\to0$ ed è diversa da zero per $h\neq0$;
per il limite notevole \eqref{eq:ln-log} e per il \cref{teo:limite-composta}, l'ultimo
fattore tende a $1$. Dunque il rapporto incrementale tende a $\frac1x$, cioè vale la
\eqref{eq:derivata-log}. Per una base qualsiasi si ha $\log_{a}x=\frac{\log x}{\log a}$
\eqref{eq:log4}, formula valida anche per esponenti reali (\cref{sec:monotone-inverse}),
e quindi $D\log_{a}x=\frac1{x\log a}$; infine $\frac1{\log a}=\log_{a}e$, sempre per la
\eqref{eq:log4}.
\end{dimostrazione}

Risulta ora chiaro l'interesse nel considerare logaritmi in base $e$: dato che
$\log_{e}e=1$, la derivata del logaritmo in base $e$ di $x$ è semplicemente $\frac1x$.

\begin{proposizione}[Derivata dell'esponenziale]\label{prop:derivata-exp}
Per ogni $x\in\R$ e ogni $a>0$ si ha
\begin{equation}\label{eq:derivata-exp}
  De^{x}=e^{x}, \qquad Da^{x}=a^{x}\log a .
\end{equation}
\end{proposizione}

\begin{schema}
\Ipotesi{$x\in\R$; \ $a>0$.}
\Tesi{$e^{x}$ e $a^{x}$ sono derivabili in $x$, \ $De^{x}=e^{x}$, \ $Da^{x}=a^{x}\log a$.}
\end{schema}

\begin{strategia}
L'esponenziale è la funzione inversa del logaritmo, di cui ora conosciamo la derivata
$\frac1x$, che non si annulla mai: il teorema di derivazione delle funzioni inverse
(\cref{teo:derivata-inversa}) dà la derivata di $e^{y}$ come reciproco di $\frac1x$,
calcolato nel punto $x=e^{y}$. Per una base qualsiasi si scrive $a^{x}=e^{x\log a}$ e si
applica la regola della catena. Dopo la dimostrazione indicheremo anche una via diretta,
basata sul limite notevole di $\frac{e^{h}-1}{h}$.
\end{strategia}

\begin{dimostrazione}
La funzione $y=\log x$ è continua, strettamente crescente e invertibile, e la sua
inversa è $x=e^{y}$. Fissato $y$, sia $x=e^{y}$: poiché $D\log x=\frac1x\neq0$, il
teorema di derivazione delle funzioni inverse (applicato al logaritmo su un intervallo
chiuso contenente $x$ al suo interno) dà
\begin{equation}\label{eq:derivata-exp-inversa}
  De^{y}=\frac1{D\log x}=\frac1{1/x}=x=e^{y}.
\end{equation}
Usando, come si è soliti fare, il simbolo $x$ per denotare la variabile indipendente,
possiamo riscrivere la formula precedente: $De^{x}=e^{x}$. Per una base $a>0$ qualsiasi
ricordiamo che $a^{x}=e^{x\log a}$ \eqref{eq:potenza-esponenziale}; facendo uso del
teorema di derivazione delle funzioni composte otteniamo
\begin{equation}\label{eq:derivata-ax}
  Da^{x}=De^{x\log a}=e^{x\log a}\,D(x\log a)=a^{x}\log a . \qedhere
\end{equation}
\end{dimostrazione}

La formula $De^{x}=e^{x}$ si può ottenere anche direttamente dalla definizione:
$\frac{e^{x+h}-e^{x}}{h}=e^{x}\cdot\frac{e^{h}-1}{h}\to e^{x}\cdot1$, per il limite
notevole \eqref{eq:ln-exp}. Essa esprime una proprietà caratteristica dell'esponenziale
naturale: \emph{il tasso di accrescimento di $e^{x}$ è uguale al suo valore}. È per
questo che la funzione $e^{x}$ descrive tutti i fenomeni in cui la rapidità di
variazione di una grandezza è proporzionale alla grandezza stessa: la crescita di una
popolazione, il decadimento radioattivo, la scarica di un condensatore, e, come vedremo
studiando le equazioni differenziali, la risposta di molti sistemi dinamici.

\subsection{Potenze con esponente reale}

Ricordiamo la proprietà $e^{\log x}=x$ \eqref{eq:log-inversa}, che è sempre utile quando
si vuole calcolare la derivata di un esponenziale o di una potenza che non rientrano nei
casi precedenti. Analogamente a quanto fatto per $a^{x}$, si calcola la derivata della
funzione potenza $x^{b}$, con esponente $b$ \emph{reale}: per $x>0$, essendo
$x^{b}=e^{b\log x}$ \eqref{eq:potenza-esponenziale},
\begin{equation}\label{eq:derivata-xb}
  Dx^{b}=De^{b\log x}=e^{b\log x}\,D(b\log x)=x^{b}\cdot\frac bx=bx^{b-1}.
\end{equation}
La formula precedente è molto utile. È utilizzata ad esempio nei casi $b=\frac12$ (in
questo caso si riottiene la \eqref{eq:derivata-radice}) e $b=-1$:
\begin{equation}\label{eq:derivata-radice-reciproco}
\begin{gathered}
  D\sqrt x=Dx^{1/2}=\frac12x^{-1/2}=\frac1{2\sqrt x},\\
  D\frac1x=Dx^{-1}=(-1)x^{-2}=-\frac1{x^{2}} .
\end{gathered}
\end{equation}
Lo stesso artificio permette di derivare le funzioni della forma $f(x)^{g(x)}$ con
$f(x)>0$, scrivendole come $e^{g(x)\log f(x)}$. Ad esempio, per $x>0$,
\begin{equation}\label{eq:derivata-x-alla-x}
  Dx^{x}=De^{x\log x}=e^{x\log x}\Bigl(\log x+x\cdot\frac1x\Bigr)=x^{x}(\log x+1).
\end{equation}

\subsection{Funzioni trigonometriche}

\begin{proposizione}[Derivate di seno, coseno e tangente]\label{prop:derivata-sen-cos}
Per ogni $x\in\R$ si ha
\begin{equation}\label{eq:derivata-sen-cos}
  D\sen x=\cos x, \qquad D\cos x=-\sen x,
\end{equation}
e per ogni $x\neq\frac\pi2+k\pi$, $k\in\Z$,
\begin{equation}\label{eq:derivata-tg}
  D\tg x=\frac1{\cos^{2}x}=1+\tg^{2}x .
\end{equation}
\end{proposizione}

\begin{schema}
\Ipotesi{$x\in\R$ (angoli in radianti); \ per la tangente, $\cos x\neq0$.}
\Tesi{$D\sen x=\cos x$; \ $D\cos x=-\sen x$; \ $D\tg x=1+\tg^{2}x$.}
\end{schema}

\begin{strategia}
Per il seno e il coseno si scrive il rapporto incrementale, si sviluppano $\sen(x+h)$ e
$\cos(x+h)$ con le formule di addizione e si raccolgono i termini in modo che $h$ compaia
soltanto nei due limiti notevoli $\frac{\sen h}{h}\to1$ e $\frac{1-\cos h}{h}\to0$; i
fattori $\sen x$ e $\cos x$, che non dipendono da $h$, escono dal limite. La tangente è il
quoziente $\frac{\sen x}{\cos x}$: basta la regola del quoziente, insieme alla relazione
fondamentale $\sen^{2}x+\cos^{2}x=1$.
\end{strategia}

\begin{dimostrazione}
Dimostriamo la prima delle \eqref{eq:derivata-sen-cos}: facciamo uso della formula di
addizione \eqref{eq:addizione-sen} e dei limiti notevoli \eqref{eq:lim-senx-x},
\eqref{eq:lim-1-cos-x}:
\begin{equation}\label{eq:derivata-sen}
\begin{split}
  \lim_{h\to0}\frac{\sen(x+h)-\sen x}{h}
  &=\lim_{h\to0}\frac{\sen x\cos h+\sen h\cos x-\sen x}{h}\\
  &=\sen x\cdot\lim_{h\to0}\frac{\cos h-1}{h}+\cos x\cdot\lim_{h\to0}\frac{\sen h}{h}\\
  &=\sen x\cdot0+\cos x\cdot1=\cos x .
\end{split}
\end{equation}
Allo stesso modo si calcola la derivata di $\cos x$, con la \eqref{eq:addizione-cos}:
\begin{equation}\label{eq:derivata-cos}
\begin{split}
  \lim_{h\to0}\frac{\cos(x+h)-\cos x}{h}
  &=\lim_{h\to0}\frac{\cos x\cos h-\sen x\sen h-\cos x}{h}\\
  &=\cos x\cdot\lim_{h\to0}\frac{\cos h-1}{h}-\sen x\cdot\lim_{h\to0}\frac{\sen h}{h}\\
  &=\cos x\cdot0-\sen x\cdot1=-\sen x .
\end{split}
\end{equation}
La derivata della funzione $\tg x$ si calcola con la regola di derivazione del quoziente,
nei punti in cui $\cos x\neq0$, usando la relazione fondamentale
\eqref{eq:relazione-fondamentale}:
\begin{equation}\label{eq:derivata-tg-calcolo}
\begin{split}
  D\tg x=D\Bigl(\frac{\sen x}{\cos x}\Bigr)&=\frac{D(\sen x)\cos x-\sen x\,D(\cos x)}{\cos^{2}x}\\
  &=\frac{\cos^{2}x+\sen^{2}x}{\cos^{2}x}=\frac1{\cos^{2}x};
\end{split}
\end{equation}
dividendo termine a termine si ottiene anche la forma $1+\tg^{2}x$.
\end{dimostrazione}

\begin{strategia}[Un'altra via]
Giusti \cite{giusti} trasforma la differenza $\sen(x+h)-\sen x$ in un prodotto con le
formule di prostaferesi: l'incremento $h$ compare allora in un solo fattore,
$\sen\frac h2$, e l'altro fattore è il coseno del punto medio. Basta così un solo limite
notevole, $\frac{\sen t}t\to1$, insieme alla continuità del coseno; il limite di
$\frac{1-\cos h}h$ non serve.
\end{strategia}

\begin{dimostrazione}[Seconda dimostrazione]
Per ogni $a$, $b$ reali, sottraendo membro a membro le formule di addizione
\eqref{eq:addizione-sen}, \eqref{eq:addizione-cos} scritte per $a+b$ e per $a-b$, si
ottengono le formule di prostaferesi (già usate nell'\cref{es:seno-lipschitz})
\begin{equation}\label{eq:prostaferesi-ab}
\begin{gathered}
  \sen(a+b)-\sen(a-b)=2\cos a\,\sen b,\\
  \cos(a+b)-\cos(a-b)=-2\sen a\,\sen b .
\end{gathered}
\end{equation}
Con $a=x+\frac h2$ e $b=\frac h2$ la prima dà, per $h\neq0$,
\[
  \frac{\sen(x+h)-\sen x}{h}=\cos\Bigl(x+\frac h2\Bigr)\cdot\frac{\sen(h/2)}{h/2}.
\]
Per $h\to0$ il primo fattore tende a $\cos x$, per la continuità del coseno
(\cref{prop:sen-cos-continue}); il secondo tende a $1$ per la \eqref{eq:lim-senx-x} e il
\cref{teo:limite-composta}, con la funzione interna $t=\frac h2$, diversa da zero per
$h\neq0$. Quindi $D\sen x=\cos x$. Allo stesso modo, per la seconda delle
\eqref{eq:prostaferesi-ab} e per la continuità del seno,
\[
  \frac{\cos(x+h)-\cos x}{h}=-\sen\Bigl(x+\frac h2\Bigr)\cdot\frac{\sen(h/2)}{h/2}
  \to-\sen x . \qedhere
\]
\end{dimostrazione}

\begin{osservazione}[La differenza centrata del seno]\label{oss:seno-centrata}
Con $a=x$ e $b=h$ la prima delle \eqref{eq:prostaferesi-ab} dà, per la differenza
centrata del \cref{sec:def-derivata},
\[
  \frac{\sen(x+h)-\sen(x-h)}{2h}=\cos x\cdot\frac{\sen h}{h}:
\]
la differenza centrata del seno è il valore esatto $\cos x$ moltiplicato per un fattore
che tende a $1$, e il suo errore vale esattamente $\cos x\,\bigl(1-\frac{\sen h}h\bigr)$
se $\cos x>0$. Per $x=1$ e $h=0{,}1$ si trova $0{,}5403\cdot0{,}0016658\approx9{,}001\cdot
10^{-4}$, il primo valore della colonna «errore centrato» della tabella del
\cref{sec:def-derivata}. Vedremo con la formula di Taylor che $1-\frac{\sen h}h$ è circa
$\frac{h^{2}}6$: per questo quell'errore si divide per $100$ quando $h$ si divide per $10$.
\end{osservazione}

\begin{attenzione}[Gli angoli vanno misurati in radianti]
Le formule \eqref{eq:derivata-sen-cos} valgono perché $\frac{\sen h}{h}\to1$, e questo
limite vale soltanto se l'angolo $h$ è misurato in radianti. Se l'angolo $\theta$ è
espresso in gradi, si ha $\sen\theta^{\circ}=\sen\frac{\pi\theta}{180}$ e, per la regola
della catena, $\frac{d}{d\theta}\sen\theta^{\circ}=\frac{\pi}{180}\cos\theta^{\circ}$: un
fattore $\frac\pi{180}\approx0{,}01745$ che nei calcoli di stabilità e di prestazioni è
fonte di errori classici (ad esempio quando si usa una pendenza della retta di portanza
$C_{L\alpha}$ espressa «per grado» in una formula scritta «per radiante»).
\end{attenzione}

Riassumiamo nella \cref{tab:derivate-elementari} le principali formule di derivazione
trovate in questo paragrafo.

\begin{table}[htb]
\centering
\renewcommand{\arraystretch}{1.45}
\begin{tabular}{l l l}
\toprule
$f(x)$ & $f'(x)$ & validità\\
\midrule
$x^{b}$ & $bx^{b-1}$ & $x>0$ ($x\in\R$ se $b\in\N$)\\
$\log x$ & $\dfrac1x$ & $x>0$\\
$\log_{a}x$ & $\dfrac1{x\log a}$ & $x>0$\\
$e^{x}$ & $e^{x}$ & $x\in\R$\\
$a^{x}$ & $a^{x}\log a$ & $x\in\R$\\
$\sen x$ & $\cos x$ & $x\in\R$\\
$\cos x$ & $-\sen x$ & $x\in\R$\\
$\tg x$ & $\dfrac1{\cos^{2}x}=1+\tg^{2}x$ & $x\neq\frac\pi2+k\pi$\\
\bottomrule
\end{tabular}
\caption{Le derivate delle funzioni elementari. Con le regole dei
\cref{sec:operazioni-derivate,sec:composte-inverse} si derivano tutte le funzioni che da
esse si ottengono con le operazioni algebriche, la composizione e l'inversione.}
\label{tab:derivate-elementari}
\end{table}

Con la tabella e le regole di calcolo il lettore può ormai derivare qualsiasi funzione
elementare. Un sistema di calcolo simbolico come \texttt{sympy} applica esattamente le
stesse regole, ed è utile per controllare i conti.

\begin{python}[Derivate calcolate con \texttt{sympy}]
import sympy as sp

x = sp.symbols("x", positive=True)
funzioni = [x**3 * sp.sin(x), sp.sin(x**2), sp.sin(x)**2, sp.tan(x),
            sp.exp(x) / x, sp.log(x**2 + 1), x**x, sp.atan(1 / x)]
for f in funzioni:
    print(f"D[{f}] = {sp.simplify(sp.diff(f, x))}")
\end{python}
\begin{risultato}
D[x**3*sin(x)] = x**2*(x*cos(x) + 3*sin(x))
D[sin(x**2)] = 2*x*cos(x**2)
D[sin(x)**2] = sin(2*x)
D[tan(x)] = cos(x)**(-2)
D[exp(x)/x] = (x - 1)*exp(x)/x**2
D[log(x**2 + 1)] = 2*x/(x**2 + 1)
D[x**x] = x**x*(log(x) + 1)
D[atan(1/x)] = -1/(x**2 + 1)
\end{risultato}

Il lettore verifichi a mano ciascun risultato. Si noti che \texttt{sympy} scrive
$\sen2x$ al posto di $2\sen x\cos x$ (formula di duplicazione) e che l'ultima riga contiene
la derivata di $\operatorname{arctg}\frac1x$, che si calcola con quella dell'arcotangente
(\cref{sec:trig-inverse}): è opposta alla derivata di $\operatorname{arctg}x$, il che non è
un caso (\cref{es:arctg-1-su-x}).

\begin{idea}[Come varia la densità dell'aria durante una salita]
Nella troposfera dell'atmosfera standard (ISA) la densità è data, in funzione della quota
(geopotenziale) $h$ in metri, da
\begin{gather*}
  \rho(h)=\rho_{0}\Bigl(1-\frac{L\,h}{T_{0}}\Bigr)^{n},\\
  \rho_{0}=1{,}225\ \text{kg/m}^{3},\quad T_{0}=288{,}15\ \text{K},\quad L=0{,}0065\ \text{K/m},
\end{gather*}
con $n=\frac{g}{RL}-1\approx4{,}2559$, dove $g=9{,}80665$~m/s$^{2}$ è l'accelerazione di
gravità standard e $R=287{,}05$~J/(kg\,K) la costante dei gas per l'aria. Per la regola
della catena e la \eqref{eq:derivata-xb},
\[
  \frac{d\rho}{dh}=\rho_{0}\,n\Bigl(1-\frac{Lh}{T_{0}}\Bigr)^{n-1}\cdot\Bigl(-\frac L{T_{0}}\Bigr),
\]
e, se il velivolo sale con velocità verticale $v_{z}=\frac{dh}{dt}$, ancora per la regola
della catena la densità varia nel tempo con tasso $\frac{d\rho}{dt}=\frac{d\rho}{dh}\,v_{z}$.
Il programma che segue calcola questi valori per $v_{z}=5$~m/s, controllando la formula
con una differenza centrata.
\end{idea}

\begin{python}[Derivata della densità ISA rispetto alla quota e al tempo]
# Densità ISA (troposfera): rho0 (1 - L h/T0)^n, n = g/(R L) - 1
rho0, T0, L, g, R = 1.225, 288.15, 0.0065, 9.80665, 287.05
n = g / (R * L) - 1
rho = lambda h: rho0 * (1 - L * h / T0) ** n
# derivata con la regola della catena: D[u^n] = n u^(n-1) u'
drho = lambda h: rho0 * n * (1 - L * h / T0) ** (n - 1) * (-L / T0)

vz = 5.0                               # velocità verticale (m/s)
print(f"n = {n:.4f}")
print("    h      rho      drho/dh     numerica     drho/dt")
for h in (0, 2000, 5000, 10000):
    num = (rho(h + 0.01) - rho(h - 0.01)) / 0.02  # diff. centrata
    print(f"{h:5d}   {rho(h):.4f}  {drho(h):.4e}  {num:.4e}  "
          f"{drho(h) * vz:.3e}")
\end{python}
\begin{risultato}
n = 4.2559
    h      rho      drho/dh     numerica     drho/dt
    0   1.2250  -1.1760e-04  -1.1760e-04  -5.880e-04
 2000   1.0065  -1.0119e-04  -1.0119e-04  -5.060e-04
 5000   0.7361  -7.9654e-05  -7.9654e-05  -3.983e-04
10000   0.4127  -5.1162e-05  -5.1162e-05  -2.558e-04
\end{risultato}

Al livello del mare la densità diminuisce di circa $1{,}2\cdot10^{-4}$~kg/m$^{3}$ per
ogni metro di quota, e salendo a $5$~m/s di circa $5{,}9\cdot10^{-4}$~kg/m$^{3}$ al
secondo; a $10\,000$~m le variazioni sono meno della metà: dalla formula si ricava
$\frac{d\rho}{dh}=-\frac{nL}{T}\,\rho$, e a quella quota la densità è circa un terzo di quella
al suolo, mentre la temperatura $T=T_{0}-Lh$ è scesa soltanto del $23\%$ circa.

%% ----------------------------------------------------------------------------
%%  §5.6  Significato geometrico della derivata. Retta tangente
%% ----------------------------------------------------------------------------
\section{Significato geometrico della derivata. Retta tangente}
\label{sec:retta-tangente}

Sia $f$ una funzione definita in un intorno di un punto $x_{0}$ e si consideri nel piano
$x$, $y$ il grafico della funzione, come nella \cref{fig:secanti}. Ci proponiamo di
determinare l'equazione della retta $r$ passante per il punto $P_{0}$ di coordinate
$(x_{0},f(x_{0}))$ e \emph{tangente} al grafico della funzione $f$. Per le circonferenze la
tangente si definisce come la retta che ha un solo punto in comune con la curva; per i
grafici di funzioni questa definizione non va bene (la retta $y=1$ ha infiniti punti in
comune con il grafico di $\sen x$, eppure è tangente in ciascuno di essi; l'asse $y$ ha un
solo punto in comune con il grafico di $x^{2}$, ma non è tangente), e occorre una
definizione basata sul concetto di limite.

Ciò che preliminarmente è più opportuno fare è determinare l'equazione di una retta $r'$
\emph{secante} il grafico della funzione $f$ nei punti $P_{0}=(x_{0},f(x_{0}))$ e
$P=(x_{0}+h,f(x_{0}+h))$, con $h\neq0$. L'equazione di una generica retta non verticale è
$y=mx+q$; determiniamo i parametri $m$, $q$ imponendo che la retta passi per i punti dati:
\begin{equation}\label{eq:sistema-secante}
  \begin{cases}
    f(x_{0})=mx_{0}+q & \text{(passaggio per } P_{0}),\\
    f(x_{0}+h)=m(x_{0}+h)+q & \text{(passaggio per } P).
  \end{cases}
\end{equation}
Abbiamo un sistema di due equazioni nelle due incognite $m$, $q$, che si può risolvere
sottraendo la prima equazione dalla seconda. Si ottiene $m=\frac{f(x_{0}+h)-f(x_{0})}{h}$,
e poi si ricava $q$ dalla prima equazione. L'equazione della retta secante risulta essere
\begin{equation}\label{eq:retta-secante}
  y=f(x_{0})+\frac{f(x_{0}+h)-f(x_{0})}{h}\,(x-x_{0}).
\end{equation}
Il coefficiente angolare della secante è dunque il rapporto incrementale di $f$ in $x_{0}$.

\begin{figure}[htb]
\centering
\begin{tikzpicture}
\begin{axis}[width=0.9\textwidth, height=6cm, xmin=0, xmax=6.3, ymin=0, ymax=5,
  axis lines=left, xtick={1.5}, xticklabels={$x_{0}$}, ytick={1.4}, yticklabels={$f(x_{0})$},
  xlabel={$x$}, ylabel={$y$}, clip=false]
  \addplot[curva, domain=0.5:6, samples=120] {1.4+1.2*(x-1.5)-0.16*(x-1.5)^2};
  \addplot[tratteggio] coordinates {(0,1.4) (1.5,1.4) (1.5,0)};
  \foreach \hh/\op/\xe in {4/0.35/6.1, 2.5/0.55/5.85, 1.2/0.8/4.95} {
    \edef\temp{\noexpand\addplot[curva2, opacity=\op, domain=0.6:\xe]
      {1.4+(1.2-0.16*\hh)*(x-1.5)};}
    \temp
    \edef\temp{\noexpand\addplot[only marks, mark=*, mark size=1.3pt, rame, opacity=\op]
      coordinates {(1.5+\hh, {1.4+1.2*\hh-0.16*\hh*\hh})};}
    \temp
  }
  \addplot[curva3, domain=0.6:4.4] {1.4+1.2*(x-1.5)};
  \addplot[only marks, mark=*, mark size=1.6pt, blu] coordinates {(1.5,1.4)};
  \node[font=\scriptsize, anchor=south east] at (axis cs:1.5,1.45) {$P_{0}$};
  \node[font=\scriptsize, verde!50!black, anchor=west] at (axis cs:4.4,4.88) {$r$};
  \node[font=\scriptsize, rame, anchor=west] at (axis cs:6.1,3.98) {$r'$};
\end{axis}
\end{tikzpicture}
\caption{Le rette secanti $r'$ per $P_{0}$ e per un secondo punto $P$ del grafico (in
rame, sempre più marcate al diminuire di $h$) tendono, quando $P$ si avvicina a $P_{0}$,
alla retta tangente $r$ (in verde), la cui pendenza è $f'(x_{0})$.}
\label{fig:secanti}
\end{figure}

L'equazione della retta tangente, quando esiste, è il limite per $h\to0$ dell'equazione
della retta secante. Si può passare al limite nella \eqref{eq:retta-secante} se e solo se
esiste finito il limite del rapporto incrementale, cioè se e solo se $f$ è derivabile in
$x_{0}$. Ciò motiva la definizione seguente.

\begin{definizione}[Retta tangente]\label{def:retta-tangente}\index{retta tangente}
Se $f$ è derivabile in $x_{0}$, si chiama \emph{retta tangente al grafico di $f$ nel
punto $(x_{0},f(x_{0}))$} la retta di equazione
\begin{equation}\label{eq:retta-tangente}
  y=f(x_{0})+f'(x_{0})(x-x_{0}).
\end{equation}
Se $f$ è continua in $x_{0}$ e il rapporto incrementale tende a $+\infty$ o a $-\infty$,
si dice che il grafico ha in $(x_{0},f(x_{0}))$ \emph{tangente verticale}, la retta
$x=x_{0}$.
\end{definizione}

Quanto stabilito fornisce il significato geometrico della derivata. Dato che
nell'equazione della retta tangente il coefficiente della $x$ è uguale a $m=f'(x_{0})$, si
dice che \emph{la derivata di una funzione $f$ in un punto $x_{0}$ è il coefficiente
angolare della retta tangente al grafico della funzione nel punto $(x_{0},f(x_{0}))$}. La
derivata è quindi una misura della \emph{pendenza} del grafico della funzione: se
$\theta\in\bigl(-\frac\pi2,\frac\pi2\bigr)$ è l'angolo che la tangente forma con l'asse $x$,
si ha $f'(x_{0})=\tg\theta$. Ad esempio, la tangente al grafico di $x^{2}$ nel punto $(1,1)$
ha coefficiente angolare $2$ ed equazione $y=1+2(x-1)=2x-1$.

\subsection{L'approssimazione lineare}

Diamo un esempio numerico di utilizzazione dell'equazione della retta tangente,
esaminando un problema di calcolo approssimato dei valori di una funzione. Normalmente non
è immediato il calcolo del valore numerico di una funzione in un punto. Ad esempio, è
facile calcolare a mente i valori numerici delle funzioni $\sqrt x$, oppure $\sen x$,
solo per particolari valori della $x$. Al contrario, è sempre elementare calcolare i
valori numerici delle funzioni $y=mx+q$, che hanno per grafico una retta. L'idea è quella
di «sostituire» una funzione data con la sua retta tangente in un punto di ascissa
$x_{0}$, con $x_{0}$ vicino al punto $x$ in cui si vuole calcolare la funzione. Dalla
\cref{fig:approssimazione-lineare} è intuitivamente chiaro che l'errore che si commette è
tanto più piccolo, quanto più $x$ è vicino all'ascissa del punto di tangenza $x_{0}$.
Cioè la quantità $f(x_{0})+f'(x_{0})(x-x_{0})$ rappresenta un'approssimazione di $f(x)$,
tanto migliore quanto più $x$ è vicino ad $x_{0}$; scriveremo
\begin{equation}\label{eq:approssimazione-lineare}
  f(x)\simeq f(x_{0})+f'(x_{0})(x-x_{0}) \qquad(\text{se } x \text{ è vicino ad } x_{0}).
\end{equation}
Il punto $x_{0}$ va scelto in modo che sia semplice calcolare $f(x_{0})$ e $f'(x_{0})$.

\begin{figure}[htb]
\centering
\begin{tikzpicture}
\begin{axis}[width=0.9\textwidth, height=5.6cm, xmin=0, xmax=6.2, ymin=0, ymax=4.3,
  axis lines=left, xtick={2.8,4.4}, xticklabels={$x_{0}$,$x$},
  ytick={2.7035,3.4023}, yticklabels={$t(x)$,$f(x)$},
  tick label style={font=\scriptsize}, xlabel={$x$}, ylabel={$y$}]
  \addplot[curva, domain=1:4.8, samples=100] {0.6+0.2*exp(0.6*x)};
  \addplot[curva3, domain=1:4.9] {1.6732+0.64387*(x-2.8)};
  \addplot[tratteggio] coordinates {(2.8,0) (2.8,1.6732)};
  \addplot[tratteggio] coordinates {(0,3.4023) (4.4,3.4023) (4.4,0)};
  \addplot[tratteggio] coordinates {(0,2.7035) (4.4,2.7035)};
  \draw[<->, rosso, thick] (axis cs:4.52,2.7035) -- node[pos=0.85, right, font=\scriptsize] {errore} (axis cs:4.52,3.4023);
  \addplot[only marks, mark=*, mark size=1.6pt, blu] coordinates {(2.8,1.6732)};
\end{axis}
\end{tikzpicture}
\caption{L'approssimazione lineare: vicino a $x_{0}$ il grafico di $f$ (blu) è ben
approssimato dalla retta tangente (verde), $t(x)=f(x_{0})+f'(x_{0})(x-x_{0})$. L'errore
$f(x)-t(x)$ tende a zero più rapidamente di $x-x_{0}$.}
\label{fig:approssimazione-lineare}
\end{figure}

Si può dare un significato rigoroso alla scrittura \eqref{eq:approssimazione-lineare}
usando i limiti. La \eqref{eq:approssimazione-lineare} significa che non solo la
differenza tra primo e secondo membro tende a zero quando $x\to x_{0}$ (questo varrebbe
per qualunque retta passante per $P_{0}$, se $f$ è continua), ma anche che tende a zero
\emph{più rapidamente} della quantità $x-x_{0}$. Conviene introdurre una notazione.

\begin{definizione}[Simbolo di Landau «o piccolo»]\label{def:o-piccolo}\index{o piccolo}
Siano $\phi$ e $\psi$ definite in un intorno di $x_{0}$ (escluso al più $x_{0}$), con
$\psi(x)\neq0$ per $x\neq x_{0}$. Si scrive $\phi(x)=o(\psi(x))$ per $x\to x_{0}$ (si
legge «$\phi$ è o piccolo di $\psi$») se
\[
  \lim_{x\to x_{0}}\frac{\phi(x)}{\psi(x)}=0 .
\]
\end{definizione}

Ad esempio $(x-x_{0})^{2}=o(x-x_{0})$ per $x\to x_{0}$, e $1-\cos x=o(x)$ per $x\to0$
\eqref{eq:lim-1-cos-x}.

\begin{proposizione}[Caratterizzazione della derivabilità]\label{prop:approssimazione-lineare}
Sia $f$ definita in un intorno di $x_{0}$. Allora $f$ è derivabile in $x_{0}$ se e solo
se esiste un numero $m\in\R$ tale che
\begin{equation}\label{eq:sviluppo-primo-ordine}
  f(x)=f(x_{0})+m(x-x_{0})+o(x-x_{0}) \qquad\text{per } x\to x_{0};
\end{equation}
in tal caso $m=f'(x_{0})$. In particolare, se $f$ è derivabile in $x_{0}$,
\begin{equation}\label{eq:errore-tangente}
  \lim_{x\to x_{0}}\frac{f(x)-[f(x_{0})+f'(x_{0})(x-x_{0})]}{x-x_{0}}=0 .
\end{equation}
\end{proposizione}

\begin{schema}
\Ipotesi{$f$ definita in un intorno di $x_{0}$.}
\Tesi{$f$ derivabile in $x_{0}$ $\sse$ esiste $m\in\R$ con
$f(x)=f(x_{0})+m(x-x_{0})+o(x-x_{0})$ per $x\to x_{0}$; \ in tal caso $m=f'(x_{0})$.}
\end{schema}

\begin{strategia}
Tutto si basa su un'identità algebrica: l'errore che si commette sostituendo $f$ con la
retta $y=f(x_{0})+m(x-x_{0})$, diviso per $x-x_{0}$, è uguale al rapporto incrementale
diminuito di $m$. Dire che l'errore è $o(x-x_{0})$ equivale quindi a dire che il rapporto
incrementale tende a $m$, cioè che $f$ è derivabile con $f'(x_{0})=m$: le due implicazioni
si leggono sulla stessa identità, nei due versi.
\end{strategia}

\begin{dimostrazione}
Per $x\neq x_{0}$ si ha l'identità
\begin{equation}\label{eq:identita-errore}
  \frac{f(x)-[f(x_{0})+m(x-x_{0})]}{x-x_{0}}=\frac{f(x)-f(x_{0})}{x-x_{0}}-m .
\end{equation}
Se $f$ è derivabile e $m=f'(x_{0})$, il secondo membro tende a $f'(x_{0})-f'(x_{0})=0$ per
la \eqref{eq:derivata-xi}: vale la \eqref{eq:sviluppo-primo-ordine}, e anche la
\eqref{eq:errore-tangente}. Viceversa, se vale la \eqref{eq:sviluppo-primo-ordine} per un
certo $m$, il primo membro della \eqref{eq:identita-errore} tende a zero, e quindi il
rapporto incrementale $\frac{f(x)-f(x_{0})}{x-x_{0}}$ tende a $m$: $f$ è derivabile e
$f'(x_{0})=m$.
\end{dimostrazione}

La proposizione dice che la retta tangente è \emph{l'unica} retta per $P_{0}$ che
approssima il grafico «al primo ordine», cioè con un errore che è $o(x-x_{0})$: qualunque
altra retta $y=f(x_{0})+m(x-x_{0})$ con $m\neq f'(x_{0})$ commette un errore il cui
rapporto con $x-x_{0}$ tende a $f'(x_{0})-m\neq0$. In questa forma la derivabilità si
estende alle funzioni di più variabili, come vedremo nel \cref{sec:differenziabilita}.
Si usa spesso anche la notazione del \emph{differenziale}: posto $dx=x-x_{0}$ e
$df=f'(x_{0})\,dx$, la \eqref{eq:sviluppo-primo-ordine} si scrive
$\Delta f=df+o(dx)$, e il differenziale $df$ è la parte lineare dell'incremento $\Delta f$.

Ad esempio, se $f(x)=\sqrt x$, la \eqref{eq:approssimazione-lineare} diventa
\begin{equation}\label{eq:approssimazione-radice}
  \sqrt x\simeq\sqrt{x_{0}}+\frac1{2\sqrt{x_{0}}}(x-x_{0}) \qquad(\text{se } x \text{ è vicino ad } x_{0}).
\end{equation}
Volendo esprimere in forma decimale $\sqrt{80}$, scegliendo $x_{0}=81$ otteniamo
\begin{equation}\label{eq:radice-80}
  \sqrt{80}\simeq9+\frac1{2\cdot9}(-1)=9-\frac1{18}=8{,}9444\ldots
\end{equation}
(il valore esatto è $\sqrt{80}=8{,}9442\ldots$). Qualcuno forse avrà pensato che il conto
è stato possibile soltanto perché $80$ è vicino al quadrato perfetto $81$. Proviamo con
$\sqrt2$: si può calcolare $\sqrt{200}$, e poi dividere il risultato per $10$. Il quadrato
più vicino a $200$ è $196=14^{2}$; si ha quindi
\begin{equation}\label{eq:radice-200}
  \sqrt{200}\simeq14+\frac1{2\cdot14}\cdot4=14+\frac17=14{,}1428\ldots
\end{equation}
Perciò $\sqrt2\simeq1{,}41428\ldots$ (il valore esatto di $\sqrt2$ è $1{,}41421\ldots$). I
conti fatti dovrebbero aver dato un'idea dell'utilità delle derivate nella tabulazione
delle funzioni elementari. Il programma che segue confronta l'approssimazione con il
valore vero in alcuni casi, e mostra che l'errore è, con buona approssimazione,
proporzionale a $(x-x_{0})^{2}$.

\begin{python}[Quanto è buona l'approssimazione lineare di $\sqrt x$?]
import math

# tangente a sqrt in x0: sqrt(x) ~ sqrt(x0) + (x - x0) / (2 sqrt(x0))
lin = lambda x, x0: math.sqrt(x0) + (x - x0) / (2 * math.sqrt(x0))
casi = [(80, 81), (200, 196), (102, 100), (110, 100), (150, 100)]
for x, x0 in casi:
    e = lin(x, x0) - math.sqrt(x)
    print(f"sqrt({x}) ~ {lin(x, x0):.6f}  vero {math.sqrt(x):.6f}  "
          f"err {e:.1e}  rapporto {e / (x - x0)**2:.5f}")
\end{python}
\begin{risultato}
sqrt(80) ~ 8.944444  vero 8.944272  err 1.7e-04  rapporto 0.00017
sqrt(200) ~ 14.142857  vero 14.142136  err 7.2e-04  rapporto 0.00005
sqrt(102) ~ 10.100000  vero 10.099505  err 5.0e-04  rapporto 0.00012
sqrt(110) ~ 10.500000  vero 10.488088  err 1.2e-02  rapporto 0.00012
sqrt(150) ~ 12.500000  vero 12.247449  err 2.5e-01  rapporto 0.00010
\end{risultato}

Con $x_{0}=100$, passando da $x=102$ a $x=110$ l'incremento $x-x_{0}$ diventa $5$ volte
più grande e l'errore circa $25$ volte più grande: è il segno che l'errore è dell'ordine
di $(x-x_{0})^{2}$, e il rapporto $\frac{\text{errore}}{(x-x_{0})^{2}}$ è quasi costante
(per $x_{0}=100$ vale circa $0{,}000125=\frac1{8\cdot100^{3/2}}$). Torneremo su questo
interessante aspetto del calcolo differenziale con la formula di Taylor, che spiegherà
il ruolo della derivata seconda in questa costante.

\begin{idea}[Sensibilità e propagazione degli errori]
L'approssimazione lineare è lo strumento con cui l'ingegnere valuta come un errore sui
dati si propaga ai risultati: se $y=f(x)$ e il dato $x$ è affetto da un piccolo errore
$\Delta x$, l'errore sul risultato è $\Delta y\simeq f'(x)\,\Delta x$. Riprendiamo
l'anemometro del \cref{cap:funzioni-continue}: la velocità si ricava dalla differenza di
pressione misurata, $V=\sqrt{2\Delta p/\rho}$, e per la \eqref{eq:derivata-radice} e la
regola della catena
\[
  \frac{dV}{d(\Delta p)}=\frac1{\sqrt{2\rho\,\Delta p}}=\frac1{\rho V}.
\]
Un errore di $5$~Pa nella misura della differenza di pressione, al livello del mare,
produce un errore di circa $5/(1{,}225\cdot20)\approx0{,}20$~m/s a $20$~m/s, ma soltanto
di circa $0{,}04$~m/s a $100$~m/s: la derivata quantifica la sensibilità che nel capitolo
precedente la sola continuità non sapeva misurare. Lo stesso principio sta alla base
della \emph{linearizzazione} in meccanica del volo: per piccole variazioni
dell'incidenza attorno a un valore di riferimento $\alpha_{0}$ si scrive
$C_{L}\simeq C_{L}(\alpha_{0})+C_{L\alpha}\,(\alpha-\alpha_{0})$, cioè si sostituisce la curva
di portanza con la sua retta tangente, e $C_{L\alpha}=\frac{dC_{L}}{d\alpha}(\alpha_{0})$ è
proprio una derivata.
\end{idea}

\subsection{Le funzioni $x^{n}\sen\frac1x$}

Utilizzando fra l'altro il significato geometrico della derivata, studiamo la continuità
e la derivabilità, nel punto $x_{0}=0$, della funzione $f_{n}$ ($n=0,1,2$) definita su $\R$
dalla formula
\begin{equation}\label{eq:fn}
  f_{n}(0)=0, \qquad f_{n}(x)=x^{n}\sen\frac1x \quad\text{se } x\neq0 .
\end{equation}
Per $n=0$ si ottiene la funzione che vale $\sen\frac1x$ per $x\neq0$ e $0$ per $x=0$; per
$n=1$ e $n=2$ le funzioni $x\sen\frac1x$ e $x^{2}\sen\frac1x$, prolungate con il valore
$0$ in $x=0$.

\begin{itemize}
  \item La funzione $f_{0}$ non è continua nel punto $x_{0}=0$, ma presenta una
  discontinuità di seconda specie: come risulta dalla \eqref{eq:sen-1-x}, non esiste il
  limite per $x\to0$ di $f_{0}(x)$. Inoltre $f_{0}$, non essendo continua in $x_{0}=0$, non
  è neanche derivabile in tale punto (\cref{teo:derivabile-continua}).
  \item La funzione $f_{1}$, prodotto del fattore infinitesimo $x$ per il fattore limitato
  $\sen\frac1x$, converge a zero per $x\to0$ (\cref{es:x-sen-1-x}); essendo $f_{1}(0)=0$,
  la funzione è continua (anche) in $x_{0}=0$. Però non risulta derivabile in tale punto,
  perché non esiste il limite del rapporto incrementale:
  \begin{equation}\label{eq:f1-rapporto}
    \lim_{h\to0}\frac{f_{1}(h)-f_{1}(0)}{h}=\lim_{h\to0}\frac{h\sen\frac1h}{h}=\lim_{h\to0}\sen\frac1h,
  \end{equation}
  che non esiste per la \eqref{eq:sen-1-x}.
  \item La funzione $f_{2}$ è derivabile (e quindi anche continua) anche per $x_{0}=0$, e
  la derivata vale $f_{2}'(0)=0$; infatti, per l'\cref{es:x-sen-1-x},
  \begin{equation}\label{eq:f2-derivata}
    f_{2}'(0)=\lim_{h\to0}\frac{f_{2}(h)-f_{2}(0)}{h}=\lim_{h\to0}h\sen\frac1h=0 .
  \end{equation}
\end{itemize}
I risultati trovati sono riassunti nella seguente tabella.
\begin{center}
\renewcommand{\arraystretch}{1.3}
\begin{tabular}{c c c}
\toprule
funzione & continua in $x_{0}=0$ & derivabile in $x_{0}=0$\\
\midrule
$f_{0}(x)$ & no & no\\
$f_{1}(x)$ & sì & no\\
$f_{2}(x)$ & sì & sì\\
\bottomrule
\end{tabular}
\end{center}

\begin{figure}[htb]
\centering
\includegraphics[width=\textwidth]{xn_sen}
\caption{I grafici di $f_{1}$ e di $f_{2}$ vicino all'origine, compresi rispettivamente tra
$y=\pm\abs x$ e tra $y=\pm x^{2}$; a destra la derivata $f_{2}'(x)$ per $x\neq0$, che oscilla
tra valori vicini a $-1$ e a $1$, mentre $f_{2}'(0)=0$ (punto blu)
(\texttt{python/fig\_xn\_sen.py}).}
\label{fig:xn-sen}
\end{figure}

Nella \cref{fig:xn-sen} sono rappresentati i grafici delle funzioni $f_{1}$ e $f_{2}$; la
continuità di $f_{1}$, $f_{2}$ nel punto $x_{0}=0$ corrisponde a grafici «vicini»
all'origine delle coordinate quando l'ascissa $x$ è «vicina» a $x_{0}=0$ (si noti che tale
proprietà non è verificata per il grafico della funzione $f_{0}$ nella
\cref{fig:sen-1-x}). Invece la derivabilità di $f_{2}$ e la non derivabilità di $f_{1}$
corrispondono al fatto che il grafico di $f_{2}$ ammette retta tangente anche nel punto
$x_{0}=0$ (dato che $f_{2}(0)=f_{2}'(0)=0$, l'equazione \eqref{eq:retta-tangente} della
retta tangente è $y=0$), mentre il grafico di $f_{1}$ non ha retta tangente nell'origine:
le rette secanti per l'origine e per il punto $(h,f_{1}(h))$ hanno coefficiente angolare
$\sen\frac1h$, che oscilla tra $-1$ e $1$ senza stabilizzarsi, e i grafici di $y=x$,
$y=-x$, con coefficienti angolari $+1$ e $-1$, danno un'idea dell'oscillazione. Tali
proprietà sono evidenziate dalle limitazioni
\begin{equation}\label{eq:f1-f2-limitazioni}
  -\abs x\le f_{1}(x)\le\abs x, \qquad -x^{2}\le f_{2}(x)\le x^{2} \qquad\perogni x\in\R .
\end{equation}

\begin{osservazione}[Una derivata che esiste ovunque ma non è continua]\label{oss:f2-non-C1}
La funzione $f_{2}$ è derivabile su tutto $\R$: per $x\neq0$, con le regole del prodotto e
della catena,
\[
  f_{2}'(x)=2x\sen\frac1x+x^{2}\cos\frac1x\cdot\Bigl(-\frac1{x^{2}}\Bigr)
  =2x\sen\frac1x-\cos\frac1x,
\]
e in $0$ vale la \eqref{eq:f2-derivata}. Ma la funzione derivata $f_{2}'$ \emph{non è
continua} in $0$: il primo addendo tende a $0$, il secondo, $\cos\frac1x$, non ha limite
per $x\to0$ (si ragiona come per $\sen\frac1x$), e quindi $f_{2}'(x)$ non ha limite per
$x\to0$ (parte destra della \cref{fig:xn-sen}). Dunque $f_{2}$ è derivabile, ma non è di
classe $C^{1}$: l'esistenza della derivata in ogni punto non garantisce la regolarità
della funzione derivata. In questo senso la $f_{2}$ è un esempio da tenere a mente ogni
volta che si è tentati di dare per scontato che una derivata sia continua.
\end{osservazione}

%% ----------------------------------------------------------------------------
%%  §5.7  Appendice: le funzioni trigonometriche inverse
%% ----------------------------------------------------------------------------
\section{Appendice: le funzioni trigonometriche inverse}
\label{sec:trig-inverse}

Le funzioni trigonometriche $\sen x$, $\cos x$, $\tg x$ non sono monotòne su tutto $\R$ (sono
periodiche) e non esistono le loro funzioni inverse su $\R$: a ogni valore del seno
corrispondono infiniti angoli. Però possiamo restringere a un intervallo limitato
l'insieme in cui prendere in considerazione tali funzioni, in modo che risultino monotòne
nell'insieme considerato. La scelta degli intervalli è convenzionale, ma universalmente
adottata (ed è quella delle funzioni di libreria di tutti i linguaggi di programmazione).

\begin{lemma}\label{lem:monotonia-trig}
La funzione $\sen x$ è strettamente crescente nell'intervallo
$\bigl[-\frac\pi2,\frac\pi2\bigr]$; la funzione $\cos x$ è strettamente decrescente
nell'intervallo $[0,\pi]$; la funzione $\tg x$ è strettamente crescente nell'intervallo
$\bigl(-\frac\pi2,\frac\pi2\bigr)$.
\end{lemma}

\begin{schema}
\Ipotesi{$x_{1}<x_{2}$, entrambi in $\bigl[-\frac\pi2,\frac\pi2\bigr]$ (seno), in $[0,\pi]$
(coseno), in $\bigl(-\frac\pi2,\frac\pi2\bigr)$ (tangente).}
\Tesi{$\sen x_{1}<\sen x_{2}$; \ $\cos x_{1}>\cos x_{2}$; \ $\tg x_{1}<\tg x_{2}$.}
\end{schema}

\begin{strategia}
Non disponiamo ancora del legame tra il segno della derivata e la monotonia (lo
stabiliremo nel prossimo capitolo): studiamo quindi direttamente il segno delle
differenze $\sen x_{2}-\sen x_{1}$, $\cos x_{2}-\cos x_{1}$, $\tg x_{2}-\tg x_{1}$. Le
formule di prostaferesi trasformano le prime due in prodotti di seni e coseni della
semisomma e della semidifferenza, che cadono in intervalli dove il segno di seno e coseno
è noto. Per la tangente si riduce allo stesso denominatore e si riconosce al numeratore il
seno di una differenza.
\end{strategia}

\begin{dimostrazione}
Dalle formule di addizione \eqref{eq:addizione-sen}, \eqref{eq:addizione-cos}, scrivendo
$x_{2}=\frac{x_{1}+x_{2}}2+\frac{x_{2}-x_{1}}2$ e $x_{1}=\frac{x_{1}+x_{2}}2-\frac{x_{2}-x_{1}}2$,
si ottengono le \emph{formule di prostaferesi}
\begin{equation}\label{eq:prostaferesi}
\begin{aligned}
  \sen x_{2}-\sen x_{1}&=2\cos\frac{x_{1}+x_{2}}2\,\sen\frac{x_{2}-x_{1}}2,\\
  \cos x_{2}-\cos x_{1}&=-2\sen\frac{x_{1}+x_{2}}2\,\sen\frac{x_{2}-x_{1}}2 .
\end{aligned}
\end{equation}
Siano $-\frac\pi2\le x_{1}<x_{2}\le\frac\pi2$. Allora $\frac{x_{1}+x_{2}}2\in
\bigl(-\frac\pi2,\frac\pi2\bigr)$, dove il coseno è positivo, e $\frac{x_{2}-x_{1}}2\in
\bigl(0,\frac\pi2\bigr]$, dove il seno è positivo; quindi $\sen x_{2}-\sen x_{1}>0$. Se invece
$0\le x_{1}<x_{2}\le\pi$, allora $\frac{x_{1}+x_{2}}2\in(0,\pi)$ e
$\frac{x_{2}-x_{1}}2\in\bigl(0,\frac\pi2\bigr]$, dove il seno è positivo; quindi
$\cos x_{2}-\cos x_{1}<0$. Infine, se $-\frac\pi2<x_{1}<x_{2}<\frac\pi2$, per la
\eqref{eq:addizione-sen}
\[
  \tg x_{2}-\tg x_{1}=\frac{\sen x_{2}\cos x_{1}-\sen x_{1}\cos x_{2}}{\cos x_{1}\cos x_{2}}
  =\frac{\sen(x_{2}-x_{1})}{\cos x_{1}\cos x_{2}}>0,
\]
perché $x_{2}-x_{1}\in(0,\pi)$ e i due coseni sono positivi.
\end{dimostrazione}

Le formule di prostaferesi sono state usate anche nell'\cref{es:seno-lipschitz} e, nella
forma \eqref{eq:prostaferesi-ab}, nella seconda dimostrazione della
\cref{prop:derivata-sen-cos}.

\subsection{L'arcoseno}

Cominciamo con la funzione $\sen x$. Per il \cref{lem:monotonia-trig} è una funzione
strettamente crescente nell'intervallo $\bigl[-\frac\pi2,\frac\pi2\bigr]$. Consideriamo quindi
$f(x)=\sen x$, con $f:\bigl[-\frac\pi2,\frac\pi2\bigr]\to[-1,1]$. La funzione $f$ è continua e
quindi assume tutti i valori compresi tra il suo minimo ($=-1$) e il suo massimo ($=1$).
Essendo strettamente monotòna, è anche invertibile (\cref{teo:invertibilita}). Pertanto
esiste la funzione inversa
\[
  f^{-1}:[-1,1]\to\Bigl[-\frac\pi2,\frac\pi2\Bigr],
\]
che viene indicata con $f^{-1}(x)=\arcsen x$ (\emph{arcoseno di $x$}). Il nome deriva dal
fatto che, se $y=\arcsen x$, vuol dire che $y$ è uguale alla misura dell'arco, o angolo,
il cui seno è $x$ ($\sen y=x$). Per il \cref{teo:continuita-inversa} l'arcoseno è
continuo, ed è strettamente crescente. Il grafico dell'arcoseno si ottiene
immediatamente dal grafico della funzione seno, per simmetria rispetto alla bisettrice
$y=x$, come nella \cref{fig:arcsen-arccos}.

\begin{figure}[htb]
\centering
\begin{tikzpicture}
\begin{axis}[name=as, width=0.5\textwidth, height=0.5\textwidth, axis lines=middle,
  xmin=-1.9, xmax=1.9, ymin=-1.9, ymax=1.9, axis equal image,
  xtick={-1.5708,-1,1,1.5708}, xticklabels={$-\frac\pi2$,$-1$,$1$,$\frac\pi2$},
  ytick={-1.5708,-1,1,1.5708}, yticklabels={$-\frac\pi2$,$-1$,$1$,$\frac\pi2$},
  tick label style={font=\tiny}, title={\scriptsize $\sen x$ e $\arcsen x$}]
  \addplot[tratteggio, domain=-1.8:1.8] {x};
  \addplot[curva2, domain=-1.5708:1.5708, samples=80] {sin(deg(x))};
  \addplot[curva, domain=-1:1, samples=200] {rad(asin(x))};
\end{axis}
\begin{axis}[at={(as.east)}, anchor=west, xshift=0.8cm, width=0.5\textwidth,
  height=0.5\textwidth, axis lines=middle, xmin=-1.4, xmax=3.5, ymin=-1.4, ymax=3.5,
  axis equal image, xtick={-1,1,1.5708,3.14159}, xticklabels={$-1$,$1$,$\frac\pi2$,$\pi$},
  ytick={-1,1,1.5708,3.14159}, yticklabels={$-1$,$1$,$\frac\pi2$,$\pi$},
  tick label style={font=\tiny}, title={\scriptsize $\cos x$ e $\arccos x$}]
  \addplot[tratteggio, domain=-1.3:3.4] {x};
  \addplot[curva2, domain=0:3.14159, samples=80] {cos(deg(x))};
  \addplot[curva, domain=-1:1, samples=200] {rad(acos(x))};
\end{axis}
\end{tikzpicture}
\caption{A sinistra il seno ristretto a $\bigl[-\frac\pi2,\frac\pi2\bigr]$ (rame) e la sua
inversa $\arcsen x$ (blu), definita su $[-1,1]$; a destra il coseno ristretto a $[0,\pi]$
e la sua inversa $\arccos x$. I grafici delle inverse si ottengono per simmetria rispetto
alla bisettrice $y=x$ (tratteggiata); nei punti $x=\pm1$ hanno tangente verticale.}
\label{fig:arcsen-arccos}
\end{figure}

La funzione arcoseno è quindi definita nell'intervallo \emph{chiuso} $[-1,1]$. Dal teorema
di derivazione delle funzioni inverse si ottiene che la funzione $\arcsen x$ è derivabile
nell'intervallo \emph{aperto} $(-1,1)$, e la derivata vale, posto $y=\arcsen x$,
\begin{equation}\label{eq:derivata-arcsen}
\begin{split}
  D\arcsen x&=\frac1{D\sen y}=\frac1{\cos y}=\frac1{\sqrt{1-\sen^{2}y}}\\
  &=\frac1{\sqrt{1-\sen^{2}(\arcsen x)}}=\frac1{\sqrt{1-x^{2}}};
\end{split}
\end{equation}
sono state utilizzate le relazioni
\begin{equation}\label{eq:relazioni-arcsen}
  \cos y=\sqrt{1-\sen^{2}y}, \qquad \sen(\arcsen x)=x,
\end{equation}
la prima delle quali vale perché $\cos y>0$ per ogni $y\in\bigl(-\frac\pi2,\frac\pi2\bigr)$
(dalla relazione fondamentale, $\cos y=\pm\sqrt{1-\sen^{2}y}$, e il segno è determinato
dalla posizione dell'angolo). La condizione $D\sen y=\cos y\neq0$, richiesta dal
\cref{teo:derivata-inversa}, è appunto verificata per $y$ interno. Nei punti $x=\pm1$, estremi del dominio, si può parlare solo di derivate
unilaterali, e in essi $\cos y=D\sen y$ si annulla ($y=\arcsen(\pm1)=\pm\frac\pi2$). Il
calcolo diretto mostra che il rapporto incrementale tende all'infinito: per $x\to1^{-}$,
posto $y=\arcsen x\to\frac\pi2^{-}$, si ha
\[
  \frac{\arcsen x-\arcsen1}{x-1}=\frac{y-\frac\pi2}{\sen y-1}
  =\frac1{\dfrac{\sen y-\sen\frac\pi2}{y-\frac\pi2}}\to+\infty,
\]
perché il denominatore tende a $D\sen\frac\pi2=\cos\frac\pi2=0$ con valori positivi
(numeratore e denominatore del rapporto $\frac{\sen y-1}{y-\pi/2}$ sono entrambi negativi);
analogamente in $x=-1$. Graficamente, ciò corrisponde al fatto che la funzione
$\arcsen x$ ha retta tangente verticale (unilaterale) se $x=\pm1$.

\subsection{L'arcocoseno}

Passiamo alla funzione $f(x)=\cos x$. Per il \cref{lem:monotonia-trig}, $f:[0,\pi]\to[-1,1]$
è continua e strettamente decrescente; quindi è invertibile in tale intervallo. La
funzione inversa $f^{-1}:[-1,1]\to[0,\pi]$ viene indicata con $f^{-1}(x)=\arccos x$
(\emph{arcocoseno di $x$}). Il grafico si ottiene dal grafico della funzione coseno, come
nella \cref{fig:arcsen-arccos}. La funzione $y=\arccos x$ è definita e continua
nell'intervallo chiuso $[-1,1]$, è derivabile nell'intervallo aperto $(-1,1)$ e la derivata
vale, posto $y=\arccos x\in(0,\pi)$, dove $\sen y>0$,
\begin{equation}\label{eq:derivata-arccos}
\begin{split}
  D\arccos x&=\frac1{D\cos y}=\frac1{-\sen y}=\frac{-1}{\sqrt{1-\cos^{2}y}}\\
  &=\frac{-1}{\sqrt{1-\cos^{2}(\arccos x)}}=\frac{-1}{\sqrt{1-x^{2}}} .
\end{split}
\end{equation}

\begin{osservazione}\label{oss:arcsen-arccos}
Le derivate di $\arcsen x$ e di $\arccos x$ sono opposte; ciò suggerisce che la loro somma
sia costante. In effetti
\begin{equation}\label{eq:arcsen-piu-arccos}
  \arcsen x+\arccos x=\frac\pi2 \qquad\perogni x\in[-1,1].
\end{equation}
Per dimostrarlo, posto $y=\arcsen x\in\bigl[-\frac\pi2,\frac\pi2\bigr]$, l'angolo
$\frac\pi2-y$ appartiene a $[0,\pi]$ e $\cos\bigl(\frac\pi2-y\bigr)=\sen y=x$ (per la
\eqref{eq:addizione-cos}); quindi, per definizione di arcocoseno,
$\arccos x=\frac\pi2-y$. (Nel prossimo capitolo vedremo che una funzione con derivata nulla
su un intervallo è costante, e questa sarà una dimostrazione alternativa.)
\end{osservazione}

\subsection{L'arcotangente}

La più usata funzione trigonometrica inversa è quella relativa alla tangente. La funzione
$f(x)=\tg x$ è continua e strettamente crescente nell'intervallo aperto
$\bigl(-\frac\pi2,\frac\pi2\bigr)$ (\cref{prop:sen-cos-continue,lem:monotonia-trig});
inoltre
\[
  \lim_{x\to\frac\pi2^{-}}\tg x=+\infty, \qquad \lim_{x\to-\frac\pi2^{+}}\tg x=-\infty,
\]
perché $\sen x\to\pm1$ mentre $\cos x\to0$ restando positivo (si veda il riquadro
«$1/0$ non è $+\infty$» del \cref{sec:forme-indeterminate}, con la funzione $\cos x$ al
posto della successione). Dunque, dato un qualsiasi $y\in\R$, esistono $a<b$ in
$\bigl(-\frac\pi2,\frac\pi2\bigr)$ con $\tg a<y<\tg b$, e per il teorema dei valori intermedi
$y$ è un valore della tangente: $\tg$ è una biiezione di $\bigl(-\frac\pi2,\frac\pi2\bigr)$ su
$\R$. La funzione inversa $f^{-1}:\R\to\bigl(-\frac\pi2,\frac\pi2\bigr)$ viene indicata con
$f^{-1}(x)=\operatorname{arctg}x$ (\emph{arcotangente di $x$}), e ha il grafico della
\cref{fig:arctg}. È continua: per ogni $y_{0}\in\R$ la restrizione della tangente a un
intervallo $[a,b]$ come sopra ha inversa continua su $[\tg a,\tg b]$
(\cref{teo:continuita-inversa}), che contiene un intorno di $y_{0}$, e basta applicare
la località del limite.

\begin{figure}[htb]
\centering
\begin{tikzpicture}
\begin{axis}[width=0.9\textwidth, height=5.4cm, axis lines=middle, xmin=-8, xmax=8,
  ymin=-2, ymax=2, xtick={-6,-4,-2,1,2,4,6}, ytick={-1.5708,0.7854,1.5708},
  yticklabels={$-\frac\pi2$,$\frac\pi4$,$\frac\pi2$}, tick label style={font=\scriptsize},
  yticklabel style={fill=white, inner sep=1pt},
  xlabel={$x$}, ylabel={$y$}]
  \addplot[tratteggio, domain=-8:8] {1.5708};
  \addplot[tratteggio, domain=-8:8] {-1.5708};
  \addplot[tratteggio] coordinates {(1,0) (1,0.7854) (0,0.7854)};
  \addplot[curva, domain=-7.8:7.8, samples=200] {rad(atan(x))};
  \node[font=\scriptsize, anchor=south east] at (axis cs:-1.5,0.4) {$y=\operatorname{arctg}x$};
\end{axis}
\end{tikzpicture}
\caption{La funzione $\operatorname{arctg}x$: è definita, continua e strettamente crescente su
tutto $\R$, con asintoti orizzontali $y=\pm\frac\pi2$; si ha $\operatorname{arctg}1=\frac\pi4$.}
\label{fig:arctg}
\end{figure}

La funzione $y=\operatorname{arctg}x$ è definita per ogni $x$ reale. Per i limiti all'infinito
si ha
\begin{equation}\label{eq:arctg-infinito}
  \lim_{x\to+\infty}\operatorname{arctg}x=\frac\pi2, \qquad \lim_{x\to-\infty}\operatorname{arctg}x=-\frac\pi2 .
\end{equation}
Infatti, fissato $\epsilon\in\bigl(0,\frac\pi2\bigr)$, per ogni $x>k=\tg\bigl(\frac\pi2-\epsilon\bigr)$
la monotonia dà $\frac\pi2-\epsilon<\operatorname{arctg}x<\frac\pi2$; il secondo limite segue
per simmetria, essendo $\operatorname{arctg}(-x)=-\operatorname{arctg}x$. È una funzione derivabile per ogni
$x\in\R$, perché $D\tg y=\frac1{\cos^{2}y}\neq0$, e la derivata vale, posto
$y=\operatorname{arctg}x$,
\begin{equation}\label{eq:derivata-arctg}
\begin{split}
  D\operatorname{arctg}x&=\frac1{D\tg y}=\frac1{1/\cos^{2}y}
  =\frac1{(\cos^{2}y+\sen^{2}y)/\cos^{2}y}\\
  &=\frac1{1+\tg^{2}y}=\frac1{1+\tg^{2}(\operatorname{arctg}x)}=\frac1{1+x^{2}} .
\end{split}
\end{equation}

Riassumiamo nella tabella seguente i valori trovati per le derivate.
\begin{center}
\renewcommand{\arraystretch}{1.7}
\begin{tabular}{l l l l}
\toprule
$f(x)$ & dominio & immagine & $f'(x)$\\
\midrule
$\arcsen x$ & $[-1,1]$ & $\bigl[-\frac\pi2,\frac\pi2\bigr]$ & $\dfrac1{\sqrt{1-x^{2}}}$ \ \ $(\abs x<1)$\\
$\arccos x$ & $[-1,1]$ & $[0,\pi]$ & $\dfrac{-1}{\sqrt{1-x^{2}}}$ \ \ $(\abs x<1)$\\
$\operatorname{arctg}x$ & $\R$ & $\bigl(-\frac\pi2,\frac\pi2\bigr)$ & $\dfrac1{1+x^{2}}$\\
\bottomrule
\end{tabular}
\end{center}

\begin{idea}[Gli angoli aerodinamici e la funzione \texttt{atan2}]
In meccanica del volo l'orientamento della velocità rispetto al velivolo si descrive con
due angoli, ricavati dalle componenti $u$ (in avanti), $v$ (laterale), $w$ (verso il basso)
della velocità in assi corpo: l'\emph{angolo d'attacco} $\alpha=\operatorname{arctg}\frac wu$ e
l'\emph{angolo di derapata} $\beta=\arcsen\frac vV$, con $V=\sqrt{u^{2}+v^{2}+w^{2}}$.
L'arcotangente restituisce valori in $\bigl(-\frac\pi2,\frac\pi2\bigr)$ e non distingue,
ad esempio, $(u,w)=(1,1)$ da $(-1,-1)$: in volo normale $u>0$ e non c'è ambiguità, ma in
manovre estreme, in una vite o nei transitori di un modello di simulazione, $u$ può
annullarsi o diventare negativo. Per questo i programmi usano la funzione a due argomenti
$\operatorname{atan2}(w,u)$, che restituisce l'angolo in $(-\pi,\pi]$ del punto $(u,w)$,
tenendo conto del quadrante (e che è definita anche per $u=0$).
\end{idea}

\begin{python}[Angolo d'attacco e angolo di derapata]
import math

# velocità in assi corpo (m/s): u avanti, v laterale, w in basso
casi = [(60.0, 0.0, 5.0), (60.0, 3.0, -4.0), (-2.0, 0.0, 10.0)]
for u, v, w in casi:
    V = math.sqrt(u*u + v*v + w*w)
    alfa_atan = math.degrees(math.atan(w / u))  # in (-90°, 90°)
    alfa = math.degrees(math.atan2(w, u))       # vede il quadrante
    beta = math.degrees(math.asin(v / V))       # derapata
    print(f"u={u:5.1f} v={v:3.1f} w={w:5.1f}: "
          f"atan {alfa_atan:7.2f}°  atan2 {alfa:7.2f}°"
          f"  beta {beta:5.2f}°")
\end{python}
\begin{risultato}
u= 60.0 v=0.0 w=  5.0: atan    4.76°  atan2    4.76°  beta  0.00°
u= 60.0 v=3.0 w= -4.0: atan   -3.81°  atan2   -3.81°  beta  2.86°
u= -2.0 v=0.0 w= 10.0: atan  -78.69°  atan2  101.31°  beta  0.00°
\end{risultato}

Nel terzo caso la velocità, in assi corpo, è diretta quasi del tutto verso il ventre del
velivolo e leggermente all'indietro, come può accadere in una scivolata di coda o in
una manovra post-stallo estrema: l'arcotangente fornisce un
angolo d'attacco di circa $-79^{\circ}$, palesemente sbagliato, mentre $\operatorname{atan2}$
restituisce il valore corretto di circa $101^{\circ}$.

\begin{notastorica}[L'arcotangente e $\pi$]
Le funzioni trigonometriche inverse sono, storicamente, più antiche della loro notazione:
tra il Tre e il Quattrocento il matematico indiano Madhava di Sangamagrama conosceva lo sviluppo in
serie
$\operatorname{arctg}x=x-\frac{x^{3}}3+\frac{x^{5}}5-\dots$, riscoperto in Europa da James Gregory
(1671) e da Leibniz (1674), che ne dedusse, con $x=1$, la celebre formula
$\frac\pi4=1-\frac13+\frac15-\frac17+\cdots$. Formule di questo tipo, combinate con identità
come quella di John Machin (1706), $\frac\pi4=4\operatorname{arctg}\frac15-\operatorname{arctg}\frac1{239}$,
permisero per due secoli di calcolare $\pi$ con centinaia di cifre. La notazione
$\sen^{-1}x$ per l'arcoseno, ancora oggi sui tasti delle calcolatrici, fu introdotta da
John Herschel nel 1813; non va confusa con $1/\sen x$, ed è per evitare l'equivoco che nei
testi italiani si preferisce $\arcsen x$.
\end{notastorica}

%% ----------------------------------------------------------------------------
%%  §5.8  Esercizi di riepilogo
%% ----------------------------------------------------------------------------
\section{Esercizi di riepilogo}
\label{sec:esercizi-cap5}

Come nei capitoli precedenti, ogni soluzione è accompagnata da una nota sulla
\emph{scelta logica} che la guida.

\needspace{8\baselineskip}
\subsection*{La definizione di derivata}

\begin{esercizio}\label{es:derivata-definizione}
Calcolare con la definizione le derivate di $\frac1x$ (per $x\neq0$) e di $\sqrt x$ (per
$x>0$).
\begin{soluzione}
Per $h\neq0$ e $x+h\neq0$,
$\frac1h\Bigl(\frac1{x+h}-\frac1x\Bigr)=\frac{x-(x+h)}{h\,x(x+h)}=-\frac1{x(x+h)}\to-\frac1{x^{2}}$.
Per la radice razionalizziamo, con $x+h>0$:
\[
  \frac{\sqrt{x+h}-\sqrt x}{h}=\frac{(x+h)-x}{h\bigl(\sqrt{x+h}+\sqrt x\bigr)}
  =\frac1{\sqrt{x+h}+\sqrt x}\to\frac1{2\sqrt x},
\]
per la continuità della radice. Si ritrovano le \eqref{eq:derivata-radice-reciproco}.
\emph{Scelta logica:} in entrambi i casi si elimina il fattore $h$ dal denominatore con
una manipolazione algebrica (denominatore comune, razionalizzazione), dopo la quale il
limite si calcola per sostituzione.
\end{soluzione}
\end{esercizio}

\begin{esercizio}\label{es:differenza-centrata}
Dimostrare che, se $f$ è derivabile in $x$, allora
$\displaystyle\lim_{h\to0}\frac{f(x+h)-f(x-h)}{2h}=f'(x)$. Mostrare con $f(x)=\abs x$ che il
limite può esistere anche se $f$ non è derivabile.
\begin{soluzione}
Si ha
$\frac{f(x+h)-f(x-h)}{2h}=\frac12\Bigl[\frac{f(x+h)-f(x)}{h}+\frac{f(x-h)-f(x)}{-h}\Bigr]$.
Il primo rapporto tende a $f'(x)$; il secondo è il rapporto incrementale con incremento
$-h$, che tende a $0$ ed è diverso da $0$, e quindi tende anch'esso a $f'(x)$
(\cref{teo:limite-composta}). Per $f(x)=\abs x$ e $x=0$ il rapporto centrato vale
$\frac{\abs h-\abs{-h}}{2h}=0$ per ogni $h$, e tende a $0$, ma $\abs x$ non è derivabile in $0$.
\emph{Scelta logica:} si spezza il rapporto simmetrico in due rapporti incrementali
ordinari; il controesempio mostra che la differenza centrata «non vede» uno spigolo
simmetrico, il che va tenuto presente nei codici numerici.
\end{soluzione}
\end{esercizio}

\begin{esercizio}
Determinare $a$ e $b$ in modo che la funzione $f(x)=x^{2}$ per $x\le1$, $f(x)=ax+b$ per
$x>1$, sia derivabile in $x=1$.
\begin{soluzione}
Serve anzitutto la continuità (\cref{teo:derivabile-continua}): $\lim_{x\to1^{+}}f(x)=a+b$
deve uguagliare $f(1)=1$. La derivata sinistra in $1$ è quella di $x^{2}$, cioè $2$; la
derivata destra, con $b=1-a$, è
$\lim_{h\to0^{+}}\frac{a(1+h)+1-a-1}{h}=a$. Deve essere $a=2$, e quindi $b=-1$: la retta
$y=2x-1$ è la tangente alla parabola nel punto $(1,1)$. \emph{Scelta logica:} prima la
continuità, poi l'uguaglianza delle derivate unilaterali. Se fosse $a+b\neq1$, il
rapporto incrementale destro $\frac{a+b-1+ah}{h}$ tenderebbe all'infinito: senza
continuità la derivata destra non esiste nemmeno.
\end{soluzione}
\end{esercizio}

\begin{esercizio}
Dimostrare che $f(x)=\abs x^{3}$ è di classe $C^{2}$ su $\R$ ma non è derivabile tre volte in
$0$; dimostrare che $f_{3}(x)=x^{3}\sen\frac1x$ (con $f_{3}(0)=0$) è di classe $C^{1}$ ma non
è derivabile due volte in $0$.
\begin{soluzione}
Si ha $f(x)=x^{2}\cdot\abs x$, e con i ragionamenti dell'\cref{es:x-modulo-x} si trova
$f'(x)=3x\abs x$ e $f''(x)=6\abs x$ per ogni $x$: $f''$ è continua ma non derivabile in $0$.
Per $f_{3}$: per $x\neq0$, $f_{3}'(x)=3x^{2}\sen\frac1x-x\cos\frac1x$, che tende a $0$; inoltre
$f_{3}'(0)=\lim_{h\to0}h^{2}\sen\frac1h=0$. Quindi $f_{3}'$ è continua. Ma
$\frac{f_{3}'(h)-f_{3}'(0)}{h}=3h\sen\frac1h-\cos\frac1h$ non ha limite per $h\to0$.
\emph{Scelta logica:} ogni fattore $x$ in più «regolarizza» di un ordine; nei punti in cui
la funzione cambia espressione si calcolano le derivate con la definizione.
\end{soluzione}
\end{esercizio}

\needspace{8\baselineskip}
\subsection*{Regole di derivazione}

\begin{esercizio}\label{es:derivate-calcolo}
Calcolare le derivate di
\begin{gather*}
  \text{(a) } x^{2}e^{-x};\qquad \text{(b) } \log(\sen x),\ x\in(0,\pi);\qquad
  \text{(c) } \sqrt{1+x^{2}};\\
  \text{(d) } e^{-x^{2}};\qquad \text{(e) } \arcsen2x;\qquad \text{(f) } x^{\sen x},\ x>0.
\end{gather*}
\begin{soluzione}
(a) $2xe^{-x}-x^{2}e^{-x}=x(2-x)e^{-x}$. (b) $\frac{\cos x}{\sen x}$. (c)
$\frac{2x}{2\sqrt{1+x^{2}}}=\frac{x}{\sqrt{1+x^{2}}}$. (d) $-2xe^{-x^{2}}$. (e)
$\frac{2}{\sqrt{1-4x^{2}}}$ per $\abs x<\frac12$. (f) $x^{\sen x}=e^{\sen x\log x}$, quindi la
derivata è $x^{\sen x}\bigl(\cos x\log x+\frac{\sen x}x\bigr)$. \emph{Scelta logica:} si
riconosce la struttura della funzione (prodotto, composta, potenza a esponente variabile)
e si applica la regola corrispondente dall'esterno verso l'interno; per gli esponenti
variabili si passa sempre alla forma $e^{g\log f}$.
\end{soluzione}
\end{esercizio}

\begin{esercizio}
Dimostrare per induzione che $D^{n}e^{ax}=a^{n}e^{ax}$ e che
$D^{n}\sen x=\sen\bigl(x+n\frac\pi2\bigr)$ per ogni $n\in\N$.
\begin{soluzione}
Per $n=1$, $De^{ax}=ae^{ax}$ (regola della catena) e
$D\sen x=\cos x=\sen\bigl(x+\frac\pi2\bigr)$. Se le formule valgono per $n$, derivando si ha
$D^{n+1}e^{ax}=a^{n}\cdot ae^{ax}$ e
$D^{n+1}\sen x=\cos\bigl(x+n\frac\pi2\bigr)=\sen\bigl(x+(n+1)\frac\pi2\bigr)$, usando
$\cos t=\sen\bigl(t+\frac\pi2\bigr)$. \emph{Scelta logica:} una formula che dipende da $n$
si dimostra per induzione, e il passo induttivo consiste nel derivare una volta in più.
\end{soluzione}
\end{esercizio}

\begin{esercizio}
Sia $f(x)=x+e^{x}$. Dimostrare che $f$ è invertibile su $\R$ e calcolare $(f^{-1})'(1)$.
\begin{soluzione}
$f$ è continua e strettamente crescente, e tende a $\pm\infty$ per $x\to\pm\infty$; quindi
è una biiezione di $\R$ su $\R$. Si ha $f(0)=1$, cioè $f^{-1}(1)=0$, e $f'(x)=1+e^{x}$, che
non si annulla mai; per il \cref{teo:derivata-inversa},
$(f^{-1})'(1)=\frac1{f'(0)}=\frac12$. \emph{Scelta logica:} non serve conoscere
esplicitamente $f^{-1}$ (che non è esprimibile con funzioni elementari): basta conoscere il
punto $x$ in cui $f(x)=1$.
\end{soluzione}
\end{esercizio}

\begin{esercizio}\label{es:arctg-1-su-x}
Dimostrare che $\operatorname{arctg}x+\operatorname{arctg}\frac1x=\frac\pi2$ per ogni $x>0$, e
verificare che la derivata del primo membro è nulla.
\begin{soluzione}
Sia $x>0$ e $\theta=\operatorname{arctg}x\in\bigl(0,\frac\pi2\bigr)$. Allora
$\frac\pi2-\theta\in\bigl(0,\frac\pi2\bigr)$ e
$\tg\bigl(\frac\pi2-\theta\bigr)=\frac{\cos\theta}{\sen\theta}=\frac1{\tg\theta}=\frac1x$; quindi
$\operatorname{arctg}\frac1x=\frac\pi2-\theta$. La derivata è
$\frac1{1+x^{2}}+\frac1{1+1/x^{2}}\cdot\bigl(-\frac1{x^{2}}\bigr)=\frac1{1+x^{2}}-\frac1{x^{2}+1}=0$, in
accordo con l'output di \texttt{sympy} del \cref{sec:derivate-elementari}. Per $x<0$ la somma
vale $-\frac\pi2$. \emph{Scelta logica:} come per l'\cref{oss:arcsen-arccos}, l'identità si
dimostra direttamente con la definizione della funzione inversa; la derivata nulla è una
verifica, che diventerà una dimostrazione con il teorema di Lagrange.
\end{soluzione}
\end{esercizio}

\needspace{8\baselineskip}
\subsection*{Retta tangente e applicazioni}

\begin{esercizio}
Scrivere l'equazione della retta tangente al grafico di $e^{x}$ nel punto di ascissa $0$, e
a quello di $\log x$ nel punto di ascissa $1$. Usarle per stimare $e^{0{,}1}$ e $\log1{,}1$.
\begin{soluzione}
Per $e^{x}$: $f(0)=f'(0)=1$, tangente $y=1+x$. Per $\log x$: $f(1)=0$, $f'(1)=1$, tangente
$y=x-1$. Le stime sono $e^{0{,}1}\simeq1{,}1$ (valore vero $1{,}10517\ldots$) e
$\log1{,}1\simeq0{,}1$ (valore vero $0{,}09531\ldots$). Le due rette sono simmetriche
rispetto alla bisettrice, come i grafici delle due funzioni inverse. \emph{Scelta logica:}
si sceglie il punto di tangenza in cui funzione e derivata sono note esattamente; le
approssimazioni coincidono con i limiti notevoli $\frac{e^{x}-1}x\to1$ e
$\frac{\log(1+x)}x\to1$.
\end{soluzione}
\end{esercizio}

\begin{esercizio}[Il radar e il tasso di variazione della distanza]
Un velivolo vola in orizzontale alla quota $H=3000$~m con velocità costante
$V=150$~m/s, e all'istante $t=0$ passa sulla verticale di un radar. Calcolare la rapidità
con cui varia la distanza radar--velivolo all'istante $t=20$~s, e il suo limite per
$t\to+\infty$.
\begin{soluzione}
La distanza è $r(t)=\sqrt{H^{2}+V^{2}t^{2}}$. Per la regola della catena,
$r'(t)=\frac{2V^{2}t}{2\sqrt{H^{2}+V^{2}t^{2}}}=\frac{V^{2}t}{r(t)}$. A $t=20$~s si ha
$Vt=3000$~m, $r=3000\sqrt2\approx4243$~m e $r'=\frac{150^{2}\cdot20}{4243}\approx106$~m/s, cioè
$V\cos45^{\circ}$: la componente della velocità lungo la congiungente. Per $t\to+\infty$,
$r'(t)=\frac{V}{\sqrt{H^{2}/(V^{2}t^{2})+1}}\to V$. \emph{Scelta logica:} è un problema di
\emph{tassi correlati}: si scrive la relazione tra le grandezze e la si deriva rispetto al
tempo con la regola della catena. Il radar Doppler misura esattamente $r'(t)$, non $V$.
\end{soluzione}
\end{esercizio}

\begin{esercizio}[Il metodo di Newton]\label{es:newton}
Per risolvere $f(x)=0$, si parte da un'approssimazione $x_{0}$ della soluzione, si sostituisce
il grafico di $f$ con la sua retta tangente nel punto $(x_{0},f(x_{0}))$ e si prende come
nuova approssimazione lo zero della tangente; poi si ripete. Ricavare la formula
dell'iterazione e applicarla all'equazione $x^{3}+x-1=0$ e all'equazione di Keplero
dell'\cref{es:kepler}.
\begin{soluzione}
La tangente in $(x_{n},f(x_{n}))$ è $y=f(x_{n})+f'(x_{n})(x-x_{n})$; se $f'(x_{n})\neq0$ il suo
zero è
\[
  x_{n+1}=x_{n}-\frac{f(x_{n})}{f'(x_{n})}.
\]
Per $f(x)=x^{3}+x-1$ si ha $f'(x)=3x^{2}+1$; per l'equazione di Keplero,
$f(E)=E-e\sen E-M$ e $f'(E)=1-e\cos E$. Il programma che segue mostra i risultati.
\emph{Scelta logica:} è l'approssimazione lineare usata «al contrario»: invece di calcolare
$f$ vicino a $x_{0}$, si cerca dove la sua approssimazione lineare si annulla. A differenza
della bisezione, il metodo non è sempre convergente; quando converge, lo fa molto
rapidamente, come si vede nell'output; la ragione si trova con la formula di Taylor
(\cref{es:newton-quadratico}).
\end{soluzione}
\end{esercizio}

\begin{python}[Il metodo di Newton]
import math

def newton(f, df, x, passi):
    """Sostituisce x con lo zero della tangente in (x, f(x))."""
    for n in range(1, passi + 1):
        x = x - f(x) / df(x)
        print(f"  x_{n} = {x:.15f}   f(x_{n}) = {f(x):+.2e}")
    return x

print("x^3 + x - 1 = 0, partendo da x_0 = 1:")
newton(lambda x: x**3 + x - 1, lambda x: 3 * x**2 + 1, 1.0, 5)
e, M = 0.5, 1.0
print("Keplero E - e sen E = M (e = 0.5, M = 1), da E_0 = M:")
newton(lambda E: E - e * math.sin(E) - M,
       lambda E: 1 - e * math.cos(E), M, 4)
\end{python}
\begin{risultato}
x^3 + x - 1 = 0, partendo da x_0 = 1:
  x_1 = 0.750000000000000   f(x_1) = +1.72e-01
  x_2 = 0.686046511627907   f(x_2) = +8.94e-03
  x_3 = 0.682339582597314   f(x_3) = +2.82e-05
  x_4 = 0.682327803946513   f(x_4) = +2.84e-10
  x_5 = 0.682327803828019   f(x_5) = +2.22e-16
Keplero E - e sen E = M (e = 0.5, M = 1), da E_0 = M:
  x_1 = 1.576469352654799   f(x_1) = +7.65e-02
  x_2 = 1.500208268606644   f(x_2) = +1.45e-03
  x_3 = 1.498701720652660   f(x_3) = +5.66e-07
  x_4 = 1.498701133517938   f(x_4) = +8.62e-14
\end{risultato}

Con cinque iterazioni il metodo di Newton raggiunge la precisione della macchina sulla
radice della cubica, e con quattro iterazioni risolve l'equazione di Keplero con un residuo
di $10^{-13}$, contro i $34$ dimezzamenti della bisezione dell'\cref{es:kepler}. Si osservi
come, a ogni passo, il numero di cifre esatte all'incirca raddoppia (il valore di $f$ passa
da $10^{-2}$ a $10^{-5}$, poi a $10^{-10}$): è la \emph{convergenza quadratica}, che si
spiega con la formula di Taylor (\cref{es:newton-quadratico}).

%% ----------------------------------------------------------------------------
%%  §5.9  Punti fondamentali del capitolo
%% ----------------------------------------------------------------------------
\section{Punti fondamentali del capitolo}
\label{sec:sintesi-cap05}

La derivata, limite del rapporto incrementale, misura la rapidità con cui varia una
grandezza e, geometricamente, la pendenza della retta tangente al grafico. Una funzione
derivabile, vicino a un punto, è approssimata al primo ordine dalla sua tangente; poche
regole di calcolo, insieme alle derivate delle funzioni elementari, permettono di
derivare qualsiasi funzione elementare.

\begin{puntichiave}
\item \textbf{Derivata.} Se $f$ è definita in un intorno di $x$, è derivabile in $x$ quando
esiste \emph{finito} $f'(x)=\lim_{h\to0}\frac{f(x+h)-f(x)}{h}$ (\cref{def:derivata}); ciò
accade se e solo se le derivate destra e sinistra esistono finite e uguali. Una
«derivata infinita» non è una derivata.
\item \textbf{Derivabilità e continuità.} Se $f$ è derivabile in $x$, è continua in $x$
(\cref{teo:derivabile-continua}); il viceversa è falso: $\abs x$ (spigolo), $\sqrt[3]x$
(tangente verticale), $\sqrt{\abs x}$ (cuspide) e $x\sen\frac1x$ sono continue in $0$ ma
non derivabili (\cref{es:modulo-non-derivabile,es:derivate-infinite}).
\item \textbf{Regole algebriche.} Se $f$ e $g$ sono derivabili in $x$, lo sono $f\pm g$,
$fg$ e, se $g(x)\neq0$, $f/g$, con $(fg)'=f'g+fg'$ e
$(f/g)'=(f'g-fg')/g^{2}$ (\cref{teo:operazioni-derivate}); la derivazione è lineare.
Errore tipico da evitare: $(fg)'\neq f'g'$.
\item \textbf{Regola della catena.} Se $g$ è derivabile in $x$ e $f$ è derivabile in
$g(x)$, allora $(f\circ g)'(x)=f'(g(x))\,g'(x)$ (\cref{teo:derivata-composta}); nel caso
generale la dimostrazione evita la divisione per un incremento $g(x+h)-g(x)$ che può
annullarsi, grazie a una funzione ausiliaria $F$ continua in $0$.
\item \textbf{Funzione inversa.} Se $f$ è continua e strettamente monotòna in $[a,b]$,
derivabile in $x\in(a,b)$ con $f'(x)\neq0$, allora $(f^{-1})'(f(x))=1/f'(x)$
(\cref{teo:derivata-inversa}); se $f'(x)=0$, l'inversa non è derivabile in $f(x)$ e il
suo grafico ha tangente verticale (\cref{oss:inversa-derivata-nulla}).
\item \textbf{Derivate elementari.} $Dx^{b}=bx^{b-1}$, $D\log x=\frac1x$, $De^{x}=e^{x}$,
$Da^{x}=a^{x}\log a$, $D\sen x=\cos x$, $D\cos x=-\sen x$, $D\tg x=1+\tg^{2}x$
(\cref{tab:derivate-elementari}); una potenza $f(x)^{g(x)}$ con $f(x)>0$ si deriva
scrivendola come $e^{g(x)\log f(x)}$.
\item \textbf{Angoli in radianti.} Le derivate di seno e coseno valgono perché
$\frac{\sen h}{h}\to1$, cioè solo in radianti; con gli angoli in gradi compare il fattore
$\frac{\pi}{180}$, causa di errori classici con una $C_{L\alpha}$ data «per grado».
\item \textbf{Funzioni trigonometriche inverse.} $\arcsen$ e $\arccos$, definite e
continue in $[-1,1]$, sono derivabili in $(-1,1)$, con derivate rispettivamente
$\frac1{\sqrt{1-x^{2}}}$ e $-\frac1{\sqrt{1-x^{2}}}$;
$\operatorname{arctg}:\R\to\bigl(-\frac\pi2,\frac\pi2\bigr)$ ha derivata $\frac1{1+x^{2}}$
(\cref{sec:trig-inverse}).
\item \textbf{Retta tangente e approssimazione lineare.} La tangente in $x_{0}$ è
$y=f(x_{0})+f'(x_{0})(x-x_{0})$ (\cref{def:retta-tangente}); $f$ è derivabile in $x_{0}$
se e solo se $f(x)=f(x_{0})+m(x-x_{0})+o(x-x_{0})$ per qualche $m$, e allora $m=f'(x_{0})$
(\cref{prop:approssimazione-lineare}): la tangente è l'unica retta che approssima il
grafico al primo ordine, e l'errore è tipicamente dell'ordine di $(x-x_{0})^{2}$.
\item \textbf{Derivate successive e classi $C^{n}$.} $f$ è di classe $C^{n}$ se $f^{(n)}$
esiste ed è continua; $x\abs x$ è di classe $C^{1}$ ma non è derivabile due volte in $0$
(\cref{es:x-modulo-x}), e $x^{2}\sen\frac1x$ è derivabile ovunque con derivata discontinua
in $0$ (\cref{oss:f2-non-C1}).
\item \textbf{Derivata numerica.} L'errore di troncamento è proporzionale a $h$ per la
differenza in avanti e a $h^{2}$ per quella centrata, mentre l'arrotondamento cresce
come $10^{-16}/h$: esiste un passo ottimale, sotto il quale ridurre $h$ peggiora il
risultato (\cref{fig:derivata-numerica}).
\item \textbf{Metodo di Newton.} Sostituendo $f$ con la tangente si ottiene l'iterazione
$x_{n+1}=x_{n}-f(x_{n})/f'(x_{n})$ (\cref{es:newton}); non sempre converge, ma quando
converge le cifre esatte all'incirca raddoppiano a ogni passo.
\item \textbf{Applicazioni aeronautiche.} Il variometro misura la derivata della quota
(\cref{sec:tasso}); la regola della catena dà $dq/dt=\rho V\,dV/dt$ per la
pressione dinamica e il tasso di variazione della densità ISA in salita
(\cref{sec:composte-inverse,sec:derivate-elementari}); la sensibilità
dell'anemometro, $dV/d(\Delta p)=1/(\rho V)$, e la linearizzazione
$C_{L}\simeq C_{L}(\alpha_{0})+C_{L\alpha}(\alpha-\alpha_{0})$ sono approssimazioni lineari
(\cref{sec:retta-tangente}).
\end{puntichiave}

\paragraph{Il filo del discorso.} Il capitolo ha costruito lo strumento, ma non ne ha
ancora sfruttato la conseguenza più importante: che cosa dice il segno della derivata
sull'andamento di una funzione? Nel \cref{cap:applicazioni-derivate} i teoremi di Fermat,
Rolle e Lagrange risponderanno a questa domanda, collegando la derivata alla monotonia, ai
massimi e minimi e alla convessità. Lo stesso capitolo presenterà il teorema di
L'Hôpital e la formula di Taylor: è questo lo strumento che giustifica perché l'errore
dell'approssimazione lineare e quello della differenza centrata sono dell'ordine del
quadrato dell'incremento, e perché il metodo di Newton converge quadraticamente.

